% Specification review / feedback for SQIsign authors
% Derived from implementing sqisign-selkie (Rust, NIST-I)
\documentclass[a4paper,11pt]{article}

\usepackage{amsfonts,amsmath,amssymb}
\usepackage[T1]{fontenc}
\usepackage{libertinus}
\usepackage[libertine]{newtxmath}
\usepackage{microtype}
\usepackage[dvipsnames]{xcolor}
\usepackage{enumitem}
\usepackage{booktabs}
\usepackage[colorlinks=true, urlcolor=BrickRed, linkcolor=BrickRed,
            citecolor=BrickRed]{hyperref}
\usepackage{xurl}
\usepackage{geometry}
\geometry{margin=28mm, marginparwidth=0pt}

\newcommand{\algo}[1]{Algorithm~#1}
\newcommand{\Fp}{\mathbb{F}_p}
\newcommand{\Fpii}{\mathbb{F}_{p^2}}

% Severity badges — inline colored small-caps
\newcommand{\critical}{\textsc{\color{BrickRed}Critical}}
\newcommand{\high}{\textsc{\color{BurntOrange}High}}
\newcommand{\medium}{\textsc{\color{RoyalBlue}Medium}}
\newcommand{\low}{\textsc{\color{OliveGreen}Low}}

% Finding header: severity + affected algorithms on one line
\newcommand{\finding}[2]{\noindent\textit{#1}\hfill#2\par\smallskip}

\title{SQIsign v2.0.1 Specification Review\\[0.3em]
  {\large Feedback from an Independent Rust Implementation}}
\author{Deirdre Connolly\\[0.2em]
  \small Selkie Cryptography\\
  \small\texttt{durumcrustulum@gmail.com}}
\date{\today\enspace{\small(living document)}}

\begin{document}
\maketitle

\begin{abstract}
Each item below describes something that the SQIsign
specification~\cite{sqisign-spec} (version 2.0.1, dated 2025-07-07) leaves unstated,
ambiguous, or implicit, and that we could only resolve by inspecting the
C reference implementation.  Implementing SQIsign from the v2.0.1 spec
alone does not currently produce an implementation that interoperates
with the C reference: the gaps below are not edge cases or
optimizations, but load-bearing conventions whose absence causes silent
KAT failures during signing or verification.  We file these as
constructive suggestions so that future implementors can achieve
interoperability from the spec alone.
\end{abstract}

\tableofcontents

% =====================================================================
\section{Severity scale}

\smallskip
\begin{tabular}{@{}lp{0.86\textwidth}@{}}
\toprule
\critical & The spec as written produces a non-interoperable
    implementation: verification rejects valid signatures, or signing
    produces unverifiable ones.  The spec must be amended. \\[0.4em]
\high & An implementor will almost certainly write a bug that passes
    unit tests but fails KAT or cross-implementation testing.
    Debugging requires reading the C reference source. \\[0.4em]
\medium & An implementor may write a bug depending on reasonable but
    wrong assumptions about conventions.  A short note eliminates the
    ambiguity. \\[0.4em]
\low & The spec is correct but incomplete.  A remark or forward pointer
    would help. \\
\bottomrule
\end{tabular}

% =====================================================================
\section{\texorpdfstring{$\Fpii$}{Fp2} square root canonicalization}
\label{sec:sqrt}

\finding{Affects \algo{8.12} (\textsc{DifferencePoint}), and
transitively every algorithm that consumes a torsion basis.}{\critical}

\algo{8.12} says ``compute $\delta = \sqrt{B_{XZ}^2 - B_{XX} B_{ZZ}}$''
without specifying which of the two square roots in $\Fpii$ to use.
The two roots $\delta$ and $-\delta$ produce two \emph{affinely distinct}
points $x_{P-Q}$.  These are not projectively equivalent; they are
genuinely different points.  The wrong choice cascades through
\textsc{Ladder3pt} (wrong kernel generator), the challenge isogeny
(different curve $E_{\text{chl}}$), and every subsequent verification
step.  In our implementation the $j$-invariant of $E_{\text{chl}}$ was
completely wrong until we matched the C reference's canonicalization.

The spec should define a canonical $\Fpii$ square root.  The C reference
uses the convention from Aardal et al.~\cite{ABCDK24}: negate the
output so that the real part (as an integer in $[0, p)$) is even; if the
real part is zero, negate so the imaginary part is even.  This should be
stated in \S8.1 (field arithmetic) as a normative requirement.  A note
of the form
\begin{quote}
  \itshape Throughout this document, ``compute $\sqrt{a}$'' for
  $a \in \Fpii$ means the unique root $r$ such that $r^2 = a$ and
  $\operatorname{Re}(r) \bmod 2 = 0$ (or, if $\operatorname{Re}(r)=0$,
  $\operatorname{Im}(r) \bmod 2 = 0$).
\end{quote}
would suffice.

% =====================================================================
\section{ProductToTheta coordinate convention}
\label{sec:product-to-theta}

\finding{Affects \algo{8.37} (\textsc{ProductToTheta}).}{\high}

The product theta coordinates are described abstractly via level-2 theta
functions.  The mapping from Montgomery projective coordinates $(X:Z)$ to
product theta coordinates is not stated explicitly.  The natural
``theta-like'' representation for Montgomery arithmetic is $(X-Z,\,X+Z)$,
which appears throughout the isogeny literature.  The C reference uses
raw $(X, Z)$ directly.  Using $(X-Z, X+Z)$ produces a completely wrong
change-of-basis matrix~$N$.

\algo{8.37} should state explicitly that the product theta coordinates
are formed as
\[
  (a, b, c, d) = (X_1 X_2,\; X_1 Z_2,\; Z_1 X_2,\; Z_1 Z_2)
\]
from Montgomery projective pairs $(X_1 : Z_1)$, $(X_2 : Z_2)$.

% =====================================================================
\section{``Squared theta'' definition}
\label{sec:squared-theta}

\finding{Affects \algo{8.29}, \algo{8.34}, \algo{8.38}.}{\high}

Several algorithms refer to ``squared theta coordinates'' without a
centralized definition.  In the C reference,
\texttt{to\_squared\_theta(P)} computes component-wise squaring
\emph{followed by a Hadamard transform}: $\tilde{P} =
\mathcal{H}(P^2)$.  The name suggests only squaring.  We omitted the
inner Hadamard in the generic isogeny evaluation, computing
$\mathcal{H}(\mathrm{inv} \cdot P^2)$ instead of the correct
$\mathcal{H}(\mathrm{inv} \cdot \mathcal{H}(P^2))$.

Add a definition early in \S8.5:
\begin{quote}
  \itshape The \textbf{squared theta} of a point $P = (a:b:c:d)$ is
  $\tilde{P} = \mathcal{H}(a^2, b^2, c^2, d^2)$, i.e.\ component-wise
  squaring followed by the Hadamard transform.
\end{quote}

% =====================================================================
\section{\textsc{GluingEval} cross-curve products are not explicit}
\label{sec:gluing-eval-cross}

\finding{Affects \algo{8.39} (\textsc{GluingEval}).}{\high}

The product-curve gluing evaluation runs \textsc{AddComponents}
on each curve component separately, producing
$(u_1, v_1, w_1)$ from curve~1 and $(u_2, v_2, w_2)$ from
curve~2, then assembles the output theta point from
products of these six values.  In the spec's pseudocode the
$U$, $V$, $W$ vectors are written as products without explicit
subscripts on every factor, leaving the reader to infer which
factors come from which curve.

The defining property of a product-curve formula is that every
multiplication has \emph{one} factor from curve~1 and \emph{one}
factor from curve~2.  An implementor reading
``$U_2 \leftarrow u_1 \cdot w_1$'' (same subscript on both
factors) cannot tell at the line whether this is a
deliberate within-curve product (rare) or a transcription error
for the cross-curve product $u_1 \cdot w_2$ (almost always
correct).  We made exactly this transcription error: the
resulting null point passed the on-curve check, ran through the
chain without raising, and the codomain failed only at
\textsc{SplittingIsomorphism} far downstream.

\paragraph{Recommendation.}
In \algo{8.39}, write every product factor with an explicit
subscript and verify the cross-curve property holds.  A short
note of the form ``\textit{every multiplication in
\textsc{AddComponents}'s output assembly takes one factor from
each component curve}'' would catch the typo class at code
review time.

% =====================================================================
\section{Projective inverse of the dual theta null point is not field
  inversion}
\label{sec:projective-inverse-gluing}

\finding{Affects \algo{8.37} (\textsc{GluingCodomain}) and
similar product-theta constructions where ``inverse'' appears
as a sub-step.}{\medium}

The gluing codomain construction (\algo{8.37}, transitively
referenced from \algo{8.39}) requires the ``inverse'' of the
dual theta null point $(\alpha, \beta, \gamma, 0)$.  In
projective theta coordinates this is computed via
cross-products --- pure multiplications producing a vector
projectively equivalent to the multiplicative inverse.  No
field inversion is involved.

The spec uses the word ``inverse'' without qualifying it, which
an implementor can reasonably read as the field-inverse
operation.  Field-inverting components in $\mathbb{F}_{p^2}$
is both numerically wrong (the projective and field inverses
agree only up to scaling that's lost between projective
classes) and expensive (each inverse is hundreds of
multiplications).

\paragraph{Recommendation.}
Replace ``inverse'' in \algo{8.37} with ``projective inverse
via cross-products'' and state the closed-form expressions
explicitly.  In general, in any projective- or theta-coordinate
algorithm, ``inverse'' should be qualified: an actual field
inversion at this layer is almost always a transcription error
from an affine reference, and is worth flagging in the spec.

% =====================================================================
\section{\texttt{hadamard\_bool} in the $(2,2)$-chain}
\label{sec:hadamard-bool}

\finding{Affects \algo{8.47}
(\textsc{Isogeny22ChainWithTorsion}).}{\high}

The algorithm describes each generic step uniformly.  The C reference
uses three formula variants controlled by boolean flags:
\begin{itemize}[nosep]
  \item Normal ($i < n{-}2$): codomain Hadamard-transformed;
        eval~$= \mathcal{H}(\mathrm{precomp} \cdot \mathcal{H}(P^2))$.
  \item Penultimate ($i = n{-}2$): codomain in dual form;
        eval~$= \mathrm{precomp} \cdot \mathcal{H}(P^2)$.
  \item Ultimate ($i = n{-}1$): extra input Hadamard, codomain dual;
        eval~$= \mathrm{precomp} \cdot \mathcal{H}(\mathcal{H}(P)^2)$.
\end{itemize}
The splitting step expects dual form.  Using the normal formula for all
steps feeds the splitter a Hadamard-transformed null point with no zero
$U_{i,j}(0)$ entry, causing verification to fail after an otherwise
correct 125-step chain.

A short remark in \algo{8.47} noting that the last two steps omit the
final Hadamard (and the ultimate step pre-Hadamards the input) would
prevent this.

% =====================================================================
\section{Torsion basis slot convention}
\label{sec:basis-swap}

\finding{Affects \algo{2.2}
(\textsc{TorsionBasisFromHint}), \algo{2.5}
(\textsc{ChangeOfBasis}), \algo{4.7}
(\textsc{ComputeChallengeIsogeny}), \algo{4.8}
(\textsc{SetChangeOfBasisMatrix}), and \algo{4.9}
(\textsc{Verify}).}{\critical}

The spec says \algo{2.2} returns the triple $(P, Q, P{-}Q)$.  The C
reference stores $P{-}Q$ in the \texttt{Q} slot and $Q$ in \texttt{PmQ},
so \textsc{Ladder3pt} computes $P + [m](P{-}Q)$, not $P + [m]Q$.  The
swap ensures the multiplied point avoids $(0,0)$ but is not documented.

\paragraph{Why this is critical, not medium.}
The shuffled storage layout is not a local
\textsc{TorsionBasisFromHint} quirk: it is the convention in which
$M_{\mathrm{chl}}$ (\algo{4.8}, signature byte 96 onwards) is encoded.
\textsc{Ladder3pt} reads the basis triple as $(P, P{-}Q, Q)$; both
signing's \textsc{SetChangeOfBasisMatrix} and \algo{4.9} line~14's
$(P, Q) \leftarrow M_{\mathrm{chl}} \cdot (P_{\mathrm{can}},
Q_{\mathrm{can}})$ run \textsc{LadderBiscalar} with that same
shuffled triple.  The matrix entries therefore express
``coordinates of the response basis in the $(P, P{-}Q)$ basis,''
not in the $(P, Q)$ basis the spec describes.

If the implementer reads the spec literally and assembles
$M_{\mathrm{chl}}$ against $(P, Q, P{-}Q)$, the bytes still parse,
verify still runs, and the resulting basis on the codomain
\emph{collapses to a singular pair} — typically $P_{\mathrm{chl}} =
Q_{\mathrm{chl}}$ bit-for-bit, since the matrix's rows project the
canonical basis onto the same point in the wrong frame.  The
$(2,2)$-chain inside \algo{4.9} either crashes (\textsc{lift\_basis}
fails, since the lift expects a non-degenerate triple) or reaches
its terminal theta null with $\#\{i, j : U_{i, j}(0) = 0\} = 0$.
Verification then rejects, and the diagnosis points everywhere except
at the basis layout.

The same convention must be used end-to-end: at
\textsc{TorsionBasisFromHint} output, at every basis manipulated in
\textsc{SplitAuxiliaryIsogeny}, \textsc{ComputeEvenNonBacktrackingResponse},
\textsc{ComputeChallengeIsogeny}, and at the change-of-basis matrix
encoding and decoding.  Mixing layouts at any one of those sites
breaks signing/verify interoperability with itself, not just with
the C reference.

\paragraph{Recommendation.}
State the layout once at the top of \S2.2 as a normative convention,
and reference it explicitly at every algorithm that manipulates a
torsion basis triple --- especially \algo{2.5}, \algo{4.7},
\algo{4.8}, and \algo{4.9}'s line 14.  An implementer who sees
``$(P, Q, P{-}Q)$'' uniformly in pseudocode has no way of knowing
that \textsc{Ladder3pt} expects the difference in the second slot.

% =====================================================================
\section{\textsc{TorsionBasisFromHint} fallback search must be bounded}
\label{sec:spec-find-uv-cap}

\finding{Affects \algo{2.2} (\textsc{TorsionBasisFromHint}), and
transitively verification (\algo{4.9}).}{\high}

\algo{2.2} reconstructs a $2$-torsion basis from a curve
coefficient~$A$ and a $7$-bit hint~$h$.  When $h = 0$ (the rare case
where the smallest admissible scalar~$n$ does not fit in $7$~bits),
the algorithm falls back to a search: starting at $n = 128$, increment
until the predicate $\bigl(n A \text{ on } E_A \text{ and not a square}\bigr)$
holds (or the dual predicate, when $A$ is a residue).  The
pseudocode reads ``while not satisfied: $n \leftarrow n + 1$,''
with no upper bound on~$n$.

\paragraph{What goes wrong.}
\textsc{TorsionBasisFromHint} is invoked from \algo{4.9}
(\textsc{Verify}) on three different curves
($E_{\mathrm{pk}}$, $E_{\mathrm{aux}}$, $E_{\mathrm{chl}}$) whose
coefficients are reconstructed directly from signature bytes.  An
adversary who controls those bytes can supply a curve coefficient~$A$
for which no $n$ in the entire search window satisfies the
predicate.  A literal implementation of the spec pseudocode loops
forever; verification never returns.

This is a denial-of-service against \texttt{verify}, reachable from
any unauthenticated input.  We discovered it via the Wycheproof
SQIsign suite, which contains a regression vector
(\texttt{HintOverflow} flag) for exactly this case.  A direct
spec-to-code translation hangs on that vector regardless of the
implementation language.

\paragraph{Structurally adversarial $A$.}
An algebraic refinement sharpens the threat.  For $A$ in the
prime subfield $\Fp \subset \Fpii$ with $A^4 - 4 A^2 + 2 = 0$
(concretely $A = \sqrt{2} \bmod p$ on NIST-I), the predicate is
\emph{structurally unsatisfiable}: an explicit calculation reduces
$V_n := (1+n^2)^2 + n^2 A^2 (A^2 - 4)$ to the perfect square
$(n^2 - 1)^2$, while $(1+n^2)$ toggles squareness independently
in~$\Fp$.  The two parities can never agree the way the predicate
needs, so the search exhausts at every~$n$.  Without a cap, an
adversary computes $\sqrt{2} \bmod p$ once (a closed-form constant
per parameter set) and crafts signatures whose
\texttt{curve\_aux} is exactly that subfield element; every call
to \textsc{Verify} on those signatures hangs forever.  The
existence of such structurally unsatisfiable $A$ means the
$K$-iteration cap is not merely an engineering convenience for
the implementer~--- it is the only termination guarantee.

\paragraph{Recommendation.}
\algo{2.2} should state explicitly that the fallback search is
bounded: ``stop after at most $K$ iterations and signal failure to
the caller; on signal, \textsc{Verify} rejects the signature.''
Any $K$ comfortably above the empirical worst-case suffices: a
Sage probe over $10^5$ random Fp$^2$-valued~$A$ found a
maximum of $37$ tries, so $K = 2^{10}$ leaves about an order of
magnitude of headroom while keeping the adversarial slowpath
below ${\sim}5$~ms (compared to $\sim 268$~ms at $K = 2^{16}$).
The spec text need not fix a specific~$K$, but should make the
cap mandatory and explicit, and should mention the structurally
unsatisfiable case so an implementer does not assume the cap is
``infallible by construction.''

\algo{4.9} (\textsc{Verify}) should be amended to consume the new
failure signal from \textsc{TorsionBasisFromHint} as a rejection.
The pattern is the same one already established in
\S\ref{sec:splitting-panic} (\textsc{GetIndexSplitting}'s
``uniqueness'' assumption): \emph{the spec's correctness
guarantees apply to honest inputs; verify must impose an
explicit cap and a rejection path on adversarial ones}.

% =====================================================================
\section{Matrix storage convention for $M_{\mathrm{chl}}$ and $M_{\mathrm{sk}}$}
\label{sec:matrix-storage}

\finding{Affects \algo{2.5}
(\textsc{ChangeOfBasis}), \algo{4.1}
(\textsc{KeyGen}: $M_{\mathrm{sk}}$), \algo{4.8}
(\textsc{SetChangeOfBasisMatrix}: $M_{\mathrm{chl}}$), and \algo{4.9}
(\textsc{Verify}: line 14 application of $M_{\mathrm{chl}}$).}{\critical}

The pseudocode of \algo{2.5} returns ``$x_1, x_2, x_3, x_4$ such that
$Q_1 = [x_1] P_1 + [x_2] P_2$ and $Q_2 = [x_3] P_1 + [x_4] P_2$.''  The
spec does not say how to lay these four scalars out as a $2 \times 2$
matrix on the wire (\S4.6 simply concatenates the four components).
Two natural layouts produce the same wire bytes only if the matrix is
written in column-major form $\bigl[\begin{smallmatrix} x_1 & x_3 \\
x_2 & x_4 \end{smallmatrix}\bigr]$, which is what the C reference does
in \texttt{change\_of\_basis\_matrix\_tate} and consumes via
\texttt{matrix\_application\_even\_basis} (column 0 yields the new
$P$).  An implementer who reads ``$Q_1 = [x_1] P_1 + [x_2] P_2$'' as a
\emph{row} ($x_1, x_2$ on row 0, $x_3, x_4$ on row 1) writes
\emph{transposed} bytes: positions for $x_2$ and $x_3$ get swapped.

\paragraph{What goes wrong.} The transposed encoding is not detectable
from a single signature because parsing succeeds, the matrix entries
are still bounded, and \algo{4.9} runs.  The defect manifests only
after \algo{4.9} line~14 applies $M_{\mathrm{chl}}$ to the canonical
basis: instead of $(P_{\mathrm{chl}}, Q_{\mathrm{chl}})$ the result is
$([x_1] P + [x_3] Q, [x_2] P + [x_4] Q)$, the application of
$M_{\mathrm{chl}}^{T}$.  The $(2,2)$-chain that follows then proceeds
on the wrong torsion basis, and verification rejects with no useful
signal.  $M_{\mathrm{sk}}$ has the same issue inside \algo{4.1}.

This is a different bug from the basis-slot swap of
\S\ref{sec:basis-swap}.  Even with the basis layout fixed,
the matrix-encoding ambiguity remains and produces
identical-looking failures.

\paragraph{Recommendation.} Either:
\begin{enumerate}
\item State \algo{4.8} explicitly produces the matrix
$\bigl[\begin{smallmatrix} x_1 & x_3 \\ x_2 & x_4 \end{smallmatrix}\bigr]$
(column-major, with column 0 the coordinates of $Q_1$ and column 1 the
coordinates of $Q_2$), and \algo{4.9} line~14 applies it by reading
column 0 to get the first new basis point; or
\item Re-state the wire encoding in \S4.6 in basis-coordinate terms:
\emph{``$M_{\mathrm{chl}}[\![0]\!]$ are the bytes of $x_1$;
$M_{\mathrm{chl}}[\![1]\!]$ of $x_3$; $M_{\mathrm{chl}}[\![2]\!]$ of $x_2$;
$M_{\mathrm{chl}}[\![3]\!]$ of $x_4$,''} so the transpose ambiguity
cannot survive a careful reading.
\end{enumerate}

% =====================================================================
\section{Spec's $2^{f-e}$ scaling in \algo{2.5} is wrong for \algo{4.8}}
\label{sec:matrix-shift}

\finding{Affects \algo{2.5} (\textsc{ChangeOfBasis}) and
\algo{4.8} (\textsc{SetChangeOfBasisMatrix}).}{\critical}

Line 4 of \algo{2.5} returns
$x_i \leftarrow 2^{f-e} \cdot \log_{\zeta}(\zeta_i)$.
The $2^{f-e}$ factor scales the $e$-bit dlog up to a $2^f$-precision
value, so the matrix can be applied to a full $2^f$-torsion basis.

But \algo{4.8} (lines 12--13) reduces both bases to order $2^e$
\emph{before} computing the matrix.  Once the bases are reduced,
the $2^{f-e}$ scaling is wrong: applying an $f$-bit entry to a
$2^e$-torsion point ignores everything above the $e$-bit window.
The wire encoding (\S4.6) stores $\lceil e/8 \rceil$ bytes per
entry; the spec's shift puts the dlog above this window, so the
encoded bytes are zero except for a few bits at the boundary.

The C reference's \texttt{\_change\_of\_basis\_matrix\_tate} omits
the $2^{f-e}$ scaling and stores raw dlogs.  This is undocumented
but necessary for the wire format and reduced-basis application
to work.

An implementer following \algo{2.5}'s step 4 literally produces
signatures whose $M_{\mathrm{chl}}$ entries are mostly zero on the
wire.  Verification accepts and runs but \algo{4.9}'s line 24
chain collapses to all-zero theta nulls; nothing flags
``$M_{\mathrm{chl}}$ entries unexpectedly small.''

\paragraph{Recommendation.}
Either fold the basis reduction into \algo{2.5}, or rewrite
\algo{2.5} step 4 conditionally: ``return $\log_{\zeta}(\zeta_i)$
when the bases are pre-reduced to order $2^e$; return
$2^{f-e} \cdot \log_{\zeta}(\zeta_i)$ when applied to a full
$2^f$-torsion basis.''  As written, the algorithm is correct for
its declared semantics (full-torsion application) but wrong for
how every protocol algorithm actually uses it.

% =====================================================================
\section{Cubical Tate third-argument sign convention in
  \algo{2.5}}
\label{sec:tate-thirdarg-sign}

\finding{Affects \algo{2.5} (\textsc{ChangeOfBasis}) step 1, and
transitively every algorithm whose change-of-basis matrix is
recovered by cubical Tate pairings (\algo{4.8}'s
$M_{\mathrm{chl}}$ inversion).}{\critical}

The cubical Tate pairing $t_{2^e}(P, Q, R)$ takes a third
projective argument~$R$ that fixes the sign branch of the
differential addition ladder.  The values
$t_{2^e}(P, Q, P{+}Q)$ and $t_{2^e}(P, Q, P{-}Q)$ are
multiplicative inverses of each other:
\[
  t_{2^e}(P, Q, P{-}Q) \;=\; t_{2^e}(P, Q, P{+}Q)^{-1}.
\]
Both satisfy bilinearity in their respective convention, both
test as $2^e$-th roots of unity, and the two conventions are
individually self-consistent.

\algo{2.5} step~1 writes $\zeta \leftarrow t_{2^e}(P_1, P_2)$
without naming the third argument's sign.  Step~3 then computes
$\zeta_i \leftarrow t_{2^e}(R_i, P_j, R_i \pm P_j)$ for each
target basis point~$R_i$ against the canonical pair $(P_1, P_2)$.

\paragraph{What goes wrong.}
An implementer who picks the SUM convention
($R = P_1 {+} P_2$) for the four cross-pairings~$\zeta_i$ but
DIFFERENCE
($R = P_1 {-} P_2$, e.g.\ because the basis triple already
carries $P_1 {-} P_2$ in the third slot per
\S\ref{sec:basis-swap}) for the base~$\zeta$ silently inverts
the relationship between~$\zeta$ and the~$\zeta_i$: the base
becomes~$T^{-1}$ where the cross-pairings remain~$T_i$, and
$\log_\zeta(\zeta_i)$ returns~$-x_i \bmod 2^e$ instead of~$x_i$.

For a target basis with all coefficients in~$\{1, 2\}$ (the
typical case in \algo{4.8}), every dlog comes out near
$2^e - 1$ or $2^e - 2$, so every matrix cell ends with its low
limbs entirely set to~$\mathtt{0xff\ldots ff}$ with at most a
few low bits differing.  Verification accepts the matrix
syntactically (entries within bounds), applies it to the
canonical basis, and obtains a \emph{non-degenerate but wrong}
basis.  The $(2,2)$-chain then runs to completion, lands on a
non-product theta null, and verify rejects.  Without
intermediate diagnostics the symptom looks like ``the chain is
just wrong'' rather than a sign-flip in the pairing setup.

The bug is invisible to most unit tests because each individual
convention is internally consistent.  Bilinearity-in-the-first
test with SUM holds.  Bilinearity-in-the-first test with
DIFFERENCE holds.  Antisymmetry holds.  Round-trip
\textsc{from\_bases} test with a known matrix is the only
shape that fires --- and even then only when the test uses
small positive matrix entries (small negatives masquerade as
near-$2^e$ values and the symptom is washed out).

\paragraph{Recommendation.}
State the third-argument convention explicitly in \algo{2.5}.
The natural choice is the SUM convention everywhere
(consistent with Jacobian addition's natural output), in which
case step~1 should read $\zeta \leftarrow t_{2^e}(P_1, P_2,
P_1{+}P_2)$.  Equivalently, an implementer who reads
\algo{8.12} (\textsc{DifferencePoint}) and naturally has
$P_1{-}P_2$ available must compute~$P_1{+}P_2$ separately for
step~1 rather than reusing the available difference.  Either
convention is mathematically fine; mixing them is catastrophic.

% =====================================================================
\section{Matrix exponent in \algo{4.8} omits \texttt{HD\_extra\_torsion}}
\label{sec:matrix-exponent}

\finding{Affects \algo{4.8} (\textsc{SetChangeOfBasisMatrix}) and
\algo{4.9} (\textsc{Verify}).}{\high}

\algo{4.8}'s pseudocode (lines 12--13) describes scaling the bases
by $2^{f - e_{\mathrm{rsp}}' - r_{\mathrm{rsp}}}$, leaving them at
order $2^{e_{\mathrm{rsp}}' + r_{\mathrm{rsp}}}$, and computing the
matrix at this exponent.  The dlog at this exponent recovers the
low $e_{\mathrm{rsp}}' + r_{\mathrm{rsp}}$ bits of each scalar
relationship; the high bits are silently zero.

But \algo{4.9}'s line 24 $(2,2)$-chain consumes
$\texttt{HD\_extra\_torsion} = 2$ bits of additional torsion: it
needs the basis at order $2^{e_{\mathrm{rsp}}' + r_{\mathrm{rsp}} + 2}$.
The C reference does both: it reduces to $2^{e_{\mathrm{rsp}}' +
r_{\mathrm{rsp}} + 2}$ \emph{and} computes the matrix at the same
exponent.  The two extra bits in the matrix entries are essential
for verify to reproduce the basis the chain needs.

The spec's \algo{4.8} omits the $+2$ silently.  An implementer who
follows the pseudocode exactly produces a matrix whose entries
truncate two high bits.  Verify then runs but its post-matrix basis
is wrong in those two bits, and the chain fails with all-zero theta
nulls — same symptom as the previous two findings, indistinguishable
from those without inspecting matrix entries directly.

\paragraph{Recommendation.}  Add ``$+ \texttt{HD\_extra\_torsion}$''
to the exponents in \algo{4.8} lines 11 and 13, matching the C
reference's \texttt{reduced\_order = pow\_dim2\_deg\_resp +
HD\_extra\_torsion + two\_resp\_length}.  This makes the dependency
on \algo{4.9}'s torsion needs explicit at the matrix-construction
site, where it can actually be implemented.

% =====================================================================
\section{Normalization of $(A_{24} : C_{24})$}
\label{sec:normalize-a24}

\finding{Affects \algo{2.2}
(\textsc{TorsionBasisFromHint}).}{\medium}

The C reference normalizes to $((A{+}2)/4 : 1)$ before basis generation.
If the curve arrives from an isogeny codomain with $C_{24} \neq 1$, the
Montgomery ladder produces a different projective representative, which
propagates through \textsc{DifferencePoint} and \textsc{Ladder3pt}.

Note in \algo{2.2} that the curve must be in normalized form before
computing the basis.

% =====================================================================
\section{Gluing requires Jacobian $y$-coordinates}
\label{sec:gluing-jacobian}

\finding{Affects \algo{8.39} (\textsc{GluingEval}).}{\low}

Montgomery $x$-only arithmetic cannot distinguish $P+Q$ from $P-Q$.  The
gluing evaluation needs this distinction.  The C reference lifts to
Jacobian $(x,y,z)$ via Okeya--Sakurai for this step---the only place in
verification that needs $y$-coordinates.  Note this in \algo{8.39}.

% =====================================================================
\section{Jacobian doubling for balanced strategy}
\label{sec:jacobian-doubling}

\finding{Affects \algo{8.47}, Phase~1.}{\low}

The spec does not specify whether Phase~1 doublings use Montgomery or
Jacobian coordinates.  The C reference doubles in Jacobian, then converts
via $\texttt{jac\_to\_xz}{:}\; (x,y,z) \mapsto (x, z^2)$.  The theta
coordinate computations are sensitive to the projective representative.
Note this, or at minimum state that the representative matters.

% =====================================================================
\section{\textsc{DifferencePoint} normalization factor}
\label{sec:difference-normalization}

\finding{Affects \algo{8.12} (\textsc{DifferencePoint}).}{\low}

The C reference multiplies $B_{XX}$, $B_{XZ}$, $B_{ZZ}$ by
$C \cdot \overline{C}^{\,2} \cdot \overline{Z_P}^{\,2} \cdot
\overline{Z_Q}^{\,2}$ (Galois conjugation) before the quadratic solve.
This makes the discriminant a fourth power in~$\Fp$, reducing the
$\Fpii$ sqrt to an $\Fp$ sqrt.  We initially used $(Z_P Z_Q)^2$, which
lacks this property since $\overline{z}^{\,2} \neq z^2$ in general.
State the normalization factor and its purpose in \algo{8.12}.

% =====================================================================

\begin{thebibliography}{9}

\bibitem{sqisign-spec}
{SQIsign team},
\textit{SQIsign: Compact Post-Quantum Signatures from Quaternions and Isogenies},
Specification version 2.0.1, dated 2025-07-07,
\url{https://sqisign.org/spec/sqisign-20250707.pdf}.

\bibitem{ABCDK24}
M.~N.~H.~Aardal, M.~P.~Azimova, E.~L.~Carri\'{o}n, L.~Dillinger, and
M.~J.~Kannwischer,
\textit{Optimized One-Dimensional SQIsign Verification on Intel and Cortex-M4},
Cryptology ePrint Archive, Paper 2024/1563,
\url{https://eprint.iacr.org/2024/1563}.

\bibitem{CHR17}
C.~Costello, H.~Hisil, and B.~Smith,
\textit{Faster Compact Diffie--Hellman: Endomorphisms on the $x$-line},
ASIACRYPT 2017,
\url{https://eprint.iacr.org/2017/518}.

\end{thebibliography}

\clearpage
\appendix

% =====================================================================
\section{Never recompute $P{-}Q$ via \textsc{DifferencePoint}}
\label{sec:pmq-recompute}

\finding{Affects \algo{4.9} (verification), lines 11--26,
and transitively the $(2,2)$-isogeny chain.}{\critical}

A torsion basis $(P,\, Q,\, P{-}Q)$ is produced by
\textsc{TorsionBasisFromHint} (\algo{2.2}).  The third component
$P{-}Q$ is computed once, via \textsc{DifferencePoint}, which
internally takes an $\Fpii$ square root.  Subsequent operations---scaling
(repeated doubling), matrix application (\algo{4.9} line~14), and the
even response isogeny (line~17)---must propagate all three points
together.  If at any point $P{-}Q$ is recomputed from $P$ and $Q$ via
\textsc{DifferencePoint}, the fresh square root may pick the opposite
branch, producing a point that is \emph{affinely distinct} from the
original $P{-}Q$.

Concretely, the two branches of \textsc{DifferencePoint} yield $x(P-Q)$
and $x(P+Q) = x(2P-Q)$; these are different points.  All subsequent
differential additions (the three-point ladder, the biladder for matrix
application, and the $(2,2)$-isogeny chain's \textsc{Ladder3pt} and
\textsc{lift\_basis}) consume $P{-}Q$ as a third input and produce wrong
results if it is silently replaced by $2P{-}Q$.

In our implementation, three separate recomputations caused failures:
\begin{enumerate}[nosep]
  \item After scaling the basis by repeated doubling, we reconstructed
        $P{-}Q$ from the doubled $P$ and $Q$ instead of doubling $P{-}Q$
        alongside them.
  \item The matrix application $M_{\mathrm{chl}} \cdot (P, Q)$ computed
        $P' - Q'$ via \textsc{DifferencePoint}$(P', Q')$ instead of a
        third biladder call $[(a{-}b)]P + [(c{-}d)]Q$.
  \item Before entering the $(2,2)$-chain, we recomputed
        $P{-}Q$ from the post-isogeny images instead of pushing the
        original $P{-}Q$ through the isogeny as a third evaluation point.
\end{enumerate}
Each of these independently caused KAT verification failure.  Fixing all
three was necessary and sufficient to make KAT vector~0 pass.

\textbf{Current spec text.} The spec does not warn against recomputing
$P{-}Q$.  Algorithms that scale or transform a basis mention only $P$
and $Q$; the third point $P{-}Q$ is not discussed.

\medskip\noindent\textbf{Recommendation.}  Add a normative note in \S2.2 or \S8.2:
\begin{quote}
  \itshape Once a basis $(P, Q, P{-}Q)$ has been produced by
  \textsc{TorsionBasisFromHint}, the third point $P{-}Q$ must be
  propagated through every subsequent operation (doubling, matrix
  application, isogeny evaluation) alongside $P$ and~$Q$.
  Recomputing $P{-}Q$ via \textsc{DifferencePoint} is not permitted:
  the square root may select a different branch, yielding an affinely
  distinct point that silently corrupts all downstream differential
  additions.
\end{quote}

% -------------------------------------------------------------------
\section{Change-of-basis matrix application convention}

\finding{Affects \algo{4.8} (\textsc{SetChangeOfBasisMatrix}),
  \algo{4.9} (verification, line~14).}{\high}

The spec does not specify whether the change-of-basis matrix
$\mathbf{M}_\mathrm{chl}$ is applied to a torsion basis $(P, Q)$
by rows or by columns.

\paragraph{What the spec says.}
\algo{4.8} line~6 writes $(P_2, Q_2) \leftarrow \mathbf{M}_1 \cdot (P_2, Q_2)$
without specifying the convention.
\algo{4.9} line~14 writes
$(P_\mathrm{chl}, Q_\mathrm{chl}) \leftarrow \mathbf{M}_\mathrm{chl} \cdot (P_\mathrm{chl}, Q_\mathrm{chl})$,
again without specifying.

\paragraph{What the C reference does.}
\texttt{matrix\_scalar\_application\_even\_basis} in \texttt{verify.c}
applies the matrix by \emph{columns}:
\begin{align*}
  R' &= [a]\,P + [c]\,Q  &&\text{(column 0: \texttt{mat[0][0]}, \texttt{mat[1][0]})} \\
  S' &= [b]\,P + [d]\,Q  &&\text{(column 1: \texttt{mat[0][1]}, \texttt{mat[1][1]})}
\end{align*}
where $\mathbf{M} = \bigl(\begin{smallmatrix} a & b \\ c & d \end{smallmatrix}\bigr)$.
An implementor who applies by rows (the natural reading of
``$\mathbf{M} \cdot (P, Q)^T$'') will produce wrong results.

\paragraph{Impact.}
Applying by rows instead of columns silently produces the wrong
torsion basis, causing the (2,2)-isogeny chain in verification to
fail (typically a panic in the splitting step).

\paragraph{Recommendation.}
State explicitly in \algo{4.8} and \algo{4.9} that the matrix is
applied by columns, or equivalently define the multiplication as
$(P', Q') = (P, Q) \cdot \mathbf{M}^T$.

% =====================================================================
\section{\textsc{RandomIdealGivenNorm}: $\beta$ sampling range is
  inclusive of~$N$}
\label{sec:random-ideal-beta-range}

\finding{Affects \algo{3.10} (\textsc{RandomIdealGivenNorm})
during signing's auxiliary-ideal construction
(\algo{4.2} line~16).}{\high}

\algo{3.10} step~11 reads
``sample $\beta = x + yi + zj + wk$ with $x, y, z, w \in [1, N]$
and $\gcd(\mathrm{nrd}(\beta), N) = 1$.''  The bracket notation
is ambiguous on the endpoint: an implementor can read
$[1, N]$ as $\{1, 2, \ldots, N\}$ (closed) or as $\{1, 2,
\ldots, N{-}1\}$ (half-open, ``$\bmod N$'' style).  Both
distributions are uniform over their respective ranges and both
preserve the algorithm's correctness as a math object.

\paragraph{Why the choice matters for byte-interop.}
On a shared DRBG the two interpretations consume the same
number of bytes per attempt (rejection sampling on a top-byte
mask), but the \emph{value} accepted differs by $1$: under
$[1, N]$ a sampled hex-bits-equal-zero case maps to $1$ (the
addition of the left endpoint), while under $[1, N{-}1]$ the
same bytes are rejected.  For every other byte pattern the two
accepted values differ by exactly~$1$.  A constant shift of
each $\beta$ coordinate by~$1$ propagates through
$\alpha = \gamma\beta$ (a quaternion product) as
$\alpha_\mathrm{ours} = \alpha_\mathrm{ref} - \gamma(1 + i + j + k)$,
which is generally not in $\mathcal{O}_0 \cdot N$.  The
resulting left ideal $\mathcal{O}_0\alpha + \mathcal{O}_0 N$
is therefore a different $\mathbb{Z}$-module on the two sides,
not just a different basis representation.

The downstream sensitivity is severe: signing's response phase
intersects this ideal with $I_{\mathrm{com,rsp}}$, then feeds
the intersection through \textsc{SuitableIdeals}.  A different
Z-module gives a different L\textsuperscript{2}-reduced basis,
a different choice of $(s, t, \beta_s, \beta_t)$, a different
$\alpha_\mathrm{rsp}$, and a different signature on the wire.

\paragraph{What the C reference does.}
\texttt{ibz\_rand\_interval(out, 1, norm)} samples uniformly on
$\{1, 2, \ldots, N\}$ inclusive of~$N$, by reading
$\lceil \log_2 N / 8 \rceil$ bytes, masking to
$\lceil \log_2(N{-}1) \rceil$ bits, rejecting samples
$> N{-}1$, and returning $\mathrm{tmp} + 1$.  This is the
closed-interval interpretation.

\paragraph{Recommendation.}
\algo{3.10} should state the convention explicitly --- the
notation should be either ``$[1, N]$ inclusive'' with a forward
pointer to ``\texttt{ibz\_rand\_interval}''-style sampling, or
``$[1, N-1]$'' (half-open) and a corresponding alternative
sampling primitive.  Either convention is fine cryptographically;
mixing them across a reference implementation and a derivative
implementation breaks KAT byte-equality and produces signatures
that fail interoperability checks.

The same ambiguity applies to every other ``sample uniformly
in $[a, b]$'' statement in the spec.  A short normative remark
near the top of \S3 stating that all such ranges are closed and
that the canonical sampler is ``$\mathrm{tmp} + a$ with
$\mathrm{tmp} \in [0, b{-}a]$ by top-byte masked rejection''
would prevent the entire class of off-by-one mismatches.

% =====================================================================
\section{\textsc{GeneralizedRepresentInteger} sampling ranges and
  divisibility}
\label{sec:represent-integer}

\finding{Affects \algo{3.12}
(\textsc{GeneralizedRepresentInteger}), and transitively
\algo{3.15} (\textsc{FixedDegreeIsogeny}) during signing.}{\medium}

\algo{3.12} describes sampling $z$ from $[1, \ldots, m]$ (line~4)
and $t$ from $[-m', \ldots, m']$ (line~5).  Two aspects of the
algorithm are easy to implement incorrectly:

\paragraph{1.\enspace Sign of $t$ matters for the isogeny condition.}
Since $M' = 4M - p(z^2 + qt^2)$ depends on~$t^2$, it is tempting
to sample only $t \geq 0$ and rely on the symmetry.  However, the
isogeny condition (line~15) checks $x - t \equiv 2 \pmod{4}$, which
\emph{does} depend on the sign of~$t$.  Restricting to $t \geq 0$
halves the effective search space for the isogeny condition, which
can cause the algorithm to exhaust its bound without finding a
solution.

\paragraph{2.\enspace Divisibility by~2 in~$\mathcal{O}$
(line~18) is not coordinate-wise.}
Line~18 says ``set~$d$ to be the biggest scalar such that
$\gamma / d \in \mathcal{O}$'' and line~19 checks $d = 2$.
For the standard order~$\mathcal{O}_0 = \mathbb{Z}[i] +
j\,\mathbb{Z}[i]$ with $q = 1$ and $\omega = i$, an element
$\gamma = x + iy + jz + kt$ satisfies $\gamma/2 \in
\mathcal{O}_0$ if and only if $x, y, z, t$ are all even.  However,
for orders with half-integer basis elements (e.g.,
$\mathcal{O}_0 = \mathbb{Z} + i\mathbb{Z} + \frac{1+j}{2}\mathbb{Z}
+ \frac{i+k}{2}\mathbb{Z}$, which is the actual $\mathcal{O}_0$
for SQIsign), divisibility by~2 in~$\mathcal{O}_0$ is
\emph{not} equivalent to all coordinates being even in the
$\{1, i, j, k\}$ basis.  The correct check must account for the
order's common denominator and basis structure.

An implementor who reads ``$\gamma / d \in \mathcal{O}$'' as
``all coordinates divisible by~$d$'' (which is correct for
$\mathbb{Z}^4$ but not for a non-trivial order) will write code
that rejects valid solutions.

\paragraph{Recommendation.}
Add a remark after line~18 of \algo{3.12} clarifying:
\begin{quote}
  \itshape The divisibility check ``$\gamma / d \in \mathcal{O}$''
  is with respect to the order~$\mathcal{O}$, not
  coordinate-wise in the $\{1, i, j, k\}$ basis.  For orders with
  a non-trivial denominator $d_{\mathcal{O}}$, the integral
  coordinates of~$\gamma$ in the order's column basis must all be
  divisible by~$d \cdot d_{\mathcal{O}}$.
\end{quote}
Also note explicitly that the range $[-m', m']$ on line~5 includes
negative values, and that both signs must be tried to satisfy the
isogeny condition on line~15.

\paragraph{3.\enspace Product order in $\gamma$ construction
(line~17).}
Line~17 writes $\gamma \leftarrow (x + \omega y + jz + \omega jt)$.
In quaternion algebra, $\omega j \neq j\omega$ in general ($ij = k$
but $ji = -k$).  The C reference's
\texttt{quat\_order\_elem\_create} computes $j \cdot \omega \cdot t$
(i.e., left-multiplying by~$j$ then by~$\omega$), which for
$\omega = i$ gives $ji \cdot t = -kt$.  Reading line~17 as
$\omega j t = ij \cdot t = kt$ produces the opposite sign on
the~$k$ component.

The difference does not affect the norm ($\mathrm{nrd}$ depends only
on $|$components$|^2$), but it \emph{does} affect the isogeny
condition (line~15) and the divisibility check (line~18), since both
depend on the parity and congruence class of the coordinates.

State whether line~17 means $(\omega j)t$ or $(j\omega)t$, or
equivalently state the product order used by
\texttt{quat\_order\_elem\_create}.

% =====================================================================
\section{Response lattice: product vs.\ intersection}
\label{sec:response-lattice}

\finding{Affects \algo{4.2} (\textsc{SQIsign.Sign}),
line~14.}{\high}

\algo{4.2} line~14 writes
\[
  \alpha_{\mathrm{rsp}} \leftarrow
  \textsc{RandomEquivalentQuaternion}(
    \overline{I_{\mathrm{com}}} \cap I_{\mathrm{sk}} \cdot I_{\mathrm{chl}})
\]
The notation ``$I_{\mathrm{sk}} \cdot I_{\mathrm{chl}}$'' is
ambiguous: it could mean the \emph{ideal product} (the $\mathbb{Z}$-span
of all $\alpha\beta$ with $\alpha \in I_{\mathrm{sk}}$,
$\beta \in I_{\mathrm{chl}}$) or the \emph{intersection}
$I_{\mathrm{sk}} \cap I_{\mathrm{chl}}$.

The C reference (\texttt{sign.c:81--85}) computes:
\begin{enumerate}[nosep]
  \item $I_{\mathrm{chl,secret}} = I'_{\mathrm{chl}} \cap
        I_{\mathrm{sk}}$ \quad (intersection, \textbf{not} product)
  \item $\overline{I_{\mathrm{com}}}$ (conjugate of $I_{\mathrm{com}}$)
  \item $\alpha_{\mathrm{rsp}} \leftarrow
        \textsc{Sample}(I_{\mathrm{chl,secret}} \cap
        \overline{I_{\mathrm{com}}})$
\end{enumerate}

Furthermore, the C ref intersects the \emph{raw} challenge ideal
$I'_{\mathrm{chl}}$ (before pushforward) with $I_{\mathrm{sk}}$,
not the pushed-forward $I_{\mathrm{chl}} = [I_{\mathrm{sk}}]_*
I'_{\mathrm{chl}}$.

An implementor who reads ``$I_{\mathrm{sk}} \cdot I_{\mathrm{chl}}$''
as an ideal product will compute a completely different lattice.  The
ideal product's norm is
$\mathrm{nrd}(I_{\mathrm{sk}}) \cdot \mathrm{nrd}(I_{\mathrm{chl}})
\approx 2^{381}$, while the intersection's norm divides
$\mathrm{lcm}(\mathrm{nrd}(I_{\mathrm{sk}}),
\mathrm{nrd}(I_{\mathrm{chl}}))$. Sampling from the wrong lattice
produces elements that generate trivial ideals (the HNF collapses to
a scaled copy of~$\mathcal{O}_0$), causing the response phase to
fail silently.

\paragraph{Recommendation.}
Replace ``$I_{\mathrm{sk}} \cdot I_{\mathrm{chl}}$'' with explicit
``$I_{\mathrm{sk}} \cap I'_{\mathrm{chl}}$'' and note that the
intersection uses the pre-pushforward challenge ideal.

% =====================================================================
\section{\textsc{LatticeSampling} algorithm is underspecified}
\label{sec:lattice-sampling}

\finding{Affects \algo{3.3}
(\textsc{LatticeSampling}), used by \algo{4.3}
(\textsc{RandomEquivalentQuaternion}) which is the inner sampling
of \algo{4.2} (\textsc{SQIsign.Sign}) line~14.}{\medium}

\algo{3.3} is described as ``sample uniformly from the lattice~$L$
elements of reduced norm at most $\rho$'' without specifying the
algorithm.  An implementor's natural reading--LLL-reduce the primal
Gram of~$L$, set per-axis bounds
$\mathrm{box}[i] = \sqrt{\rho / G_{\mathrm{red}}[i][i]}$, and reject
sample--is correct but produces a bounding parallelogram much looser
than the input ellipsoid, with a typical acceptance rate around
$10^{-4}$ at the radii used in signing.

The C reference's \texttt{quat\_lattice\_sample\_from\_ball}
(\texttt{lat\_ball.c:58}) implements a different, tighter algorithm:
LLL-reduce the \emph{dual} Gram $\mathrm{adj}(G)$ while tracking the
unimodular transformation~$U$, set
$\mathrm{box}[i] = \sqrt{\mathrm{dual}G_{\mathrm{red}}[i][i]
\cdot \rho / \det G}$, sample $y$ in the box, map
$x = U^{-\top} y$, and check $x^\top G x \le \rho$ on the original
primal Gram.  This bound corresponds to the ellipsoid's inertia
tensor (via the dual lattice's reduction) and gives an acceptance
rate within a small Hermite-constant factor of optimal.

The two algorithms produce different distributions of accepted~$\alpha$,
and downstream Kani-kernel acceptance in
\algo{4.5} (\textsc{SplitAuxiliaryIsogeny}) tracks the
$\alpha$~distribution.  An implementor following the spec literally
will produce a different per-iter retry profile than the C ref and
may attribute the timing or convergence differences to other
algorithms entirely.  We did, until cross-checking the C ref's
sampling code revealed the gap.

\paragraph{Recommendation.}
Spell out the dual-LLL procedure (and the
$\mathrm{box}[i] = \sqrt{\mathrm{dual}G_{\mathrm{red}}[i][i]
\cdot \rho / \det G}$ bound, plus the $U^{-\top}$ map) in
\algo{3.3}, or at least name it as the recommended algorithm
with a citation to a standard reference.  As stated, the spec
permits any sampling that meets the distribution, but the C
reference's choice is the load-bearing one for matching its
benchmark timings.

% =====================================================================
\section{\texorpdfstring{$L^2$}{L2}-deep-insertion iteration order
underspecified}
\label{sec:l2-deep-insertion-order}

\finding{Affects \algo{3.6} ($L^2$-DeepInsertion), used by every
LLL reduction in keygen and signing through
\texttt{quat\_lideal\_prime\_norm\_reduced\_equivalent}.}{\high}

\algo{3.6} defines the deep-insertion point as ``the smallest~$j$
such that $\bar\delta \cdot r[j][j] > t[j]$,'' i.e.\ the smallest
level at which the $L^2$-condition fails on the partial squared-norm
sequence $t[0], \ldots, t[\kappa-1]$.

The C~reference's \texttt{quat\_lll\_core}
(\texttt{quaternion/ref/generic/lll/l2.c:110-114}) uses a different
iteration order:
\begin{verbatim}
for (swap = kappa; swap > 0; swap--) {
    if (delta_bar * r[swap-1][swap-1] < lovasz[swap-1])
        break;
}
\end{verbatim}
This iterates~$s$ \emph{from~$\kappa$ down to~$1$} and breaks on the
first level where the $L^2$-condition still holds.

When the failure pattern is monotonic---fails for $j < a$, holds
for $j \ge a$---both formulations return $s = a$.  But the partial
squared norms $t[j]$ are computed from a running Gram--Schmidt
orthogonalization update on floating-point scalars (\texttt{dpe\_t}
in the C~ref).  Under rounding the failure pattern can be
\emph{non-monotonic}: failures at, say, $j = 0$ and $j = 2$ but a
hold at $j = 1$.  The spec-literal form returns $s = 0$
(smallest failing); the C~ref returns $s = 2$ (largest level above
the deepest contiguous hold).  Both are valid $L^2$-reduced bases,
but the resulting bases differ in the sign of column zero,
propagating to off-diagonal sign flips in the post-LLL Gram matrix
($G[0][2]$, $G[0][3]$ and their transposes; diagonals match
byte-for-byte).

The downstream call to
\texttt{quat\_lideal\_prime\_norm\_reduced\_equivalent}
(\texttt{lll\_applications.c}) rejection-samples short vectors~$c$
and accepts the first~$c$ whose class quadratic form
$c^\top G_{\mathrm{class}} c$ is prime.  A different post-LLL
$G_{\mathrm{class}}$ produces a different acceptance, hence a
different equivalent prime ideal, hence a different chain codomain,
hence a different public key.

\paragraph{Empirical impact.}
For~16 of~100 NIST-I keygen KATs every input that hit the
non-monotonic condition---this single iteration-order divergence
cascaded all the way to a public-key mismatch.  Switching from the
spec-literal order (smallest failing $j$, ascending) to the C~ref
order (deepest contiguous failure, descending) moved our keygen
byte-equality count from~83/100 to~99/100 in a single commit.

\paragraph{Current spec text.}
\algo{3.6} (and~\algo{3.3} which calls it) describe the
deep-insertion choice mathematically but do not specify the
iteration direction.  Both orders satisfy the spec's English
description under the implicit monotonicity assumption; floating-
point rounding can break that assumption silently.

\paragraph{Recommendation.}
Either (a)~specify the iteration order in \algo{3.6}---``iterate
$s = \kappa, \kappa-1, \ldots, 1$ and stop at the first $s$ such
that $\bar\delta \cdot r[s-1][s-1] \ge t[s-1]$''---or (b)~note that
the choice of order is implementation-defined but that KAT
byte-equality requires matching the reference's order exactly.
Without this clarification, an implementor who follows the spec's
``smallest~$j$ that fails'' will produce a non-interoperable build
and need to read \texttt{quat\_lll\_core} source to find the gap.

% =====================================================================
\section{\textsc{FixedDegreeIsogeny} missing $u^{-1}$ normalization}
\label{sec:u-inv}

\finding{Affects \algo{3.15}
(\textsc{FixedDegreeIsogeny}), and transitively \algo{3.13}
(\textsc{IdealToIsogeny}) during signing.}{\critical}

\algo{3.15} line~2 computes
$\theta \leftarrow \textsc{RepresentInteger}(u \cdot (2^{e} - u),
\mathcal{O}_t, \mathrm{true})$, and line~4 constructs the
$(2,2)$-isogeny kernel from $([u]P,\, \theta(P))$ and
$([u]Q,\, \theta(Q))$.  Since $\mathrm{nrd}(\theta) = u \cdot
(2^{e} - u)$, the endomorphism $\theta$ has degree
$u (2^{e} - u)$; the kernel we want is that of the degree-$(2^{e}-u)$
isogeny, generated by $(P,\, (\theta/u)(P))$.

The C reference multiplies all four quaternion coordinates of~$\theta$
by $u^{-1} \bmod 2^{e+2}$ before applying the endomorphism to the
torsion basis (\texttt{dim2id2iso.c}, lines 133--143).  Without this
normalization, the \textsc{ActionByTranslation} determinant of the
gluing kernel is degenerate: the splitting step finds zero (or ten)
vanishing~$U_{i,j}(0)$ indices instead of exactly one, and the
$(2,2)$-chain fails.

\paragraph{Current spec text.}
\algo{3.15} line~4 writes ``$K_1 = (u \cdot P_t,\, \theta(P_t))$''
without mentioning the $u^{-1}$ normalization.  The pseudocode is
algebraically suggestive (the $u$-torsion is killed after $f{-}2{-}e$
doublings), but in practice the non-normalized kernel is numerically
degenerate and no $(2,2)$-chain succeeds.

\paragraph{Recommendation.}
Replace line~4 of \algo{3.15} with
$(P_t,\, (\theta \cdot u^{-1} \bmod 2^{e+2})(P_t))$ and note
that $u$ is odd (hence invertible mod~$2^{e+2}$) and that the
extra~$+2$ in the modulus accounts for the gluing's order-4
structure.

% =====================================================================
\section{Fixed-precision Hermite Normal Form}
\label{sec:hnf}

\finding{Affects \algo{3.2} (\textsc{HNF}), and transitively every
algorithm that constructs or reduces an ideal lattice --- including
\algo{3.9} (\textsc{RandomEquivalentPrimeIdeal}),
\algo{3.10} (\textsc{RandomIdealGivenNorm}), and
\algo{3.15} (\textsc{FixedDegreeIsogeny}).}{\high}

\algo{3.2} describes the column-style Hermite Normal Form via
extended-GCD elimination.  The pseudocode operates over
$\mathbb{Z}$, i.e., arbitrary-precision integers.

\paragraph{What goes wrong.}
Classical HNF's intermediate entries can grow exponentially in
the rank.  For rank-4 quaternion ideal lattices at NIST-I the
initial column entries can reach ${\sim}2^{770}$~bits (the
$p \cdot g_i$ products from \algo{3.10}).  After a few xgcd
elimination steps, intermediate column entries exceed any
reasonable fixed-width budget.  At our storage width of 1920~bits
($30 \times 64$), the resulting ``HNF'' silently wraps: we
observed basis columns collapse to elements of the ambient
order~$O_0$ (reduced norms of~4 and~${\sim}2^{251}$) rather
than the intended ideal~$I$.  Every subsequent operation ---
\textsc{RandomEquivalentPrimeIdeal}, \textsc{IdealToIsogeny},
the entire response phase --- then operates on a lattice unrelated
to~$I$.

\paragraph{What the C reference does.}
The C reference uses GMP (\texttt{ibz\_t}), so coefficient growth
is absorbed by arbitrary-precision arithmetic.  The classical
algorithm works correctly there.

\paragraph{Fix.}
Use the standard Domich--Kannan--Trotter modular HNF (Cohen,
\emph{A Course in Computational Algebraic Number Theory},
\S2.4.2): given a positive integer~$D$ with $D \cdot e_i \in L$
for every standard basis vector~$e_i$, reduce every intermediate
entry modulo~$D$ after every arithmetic update.  This bounds
entries by~$D$ instead of by the exponential classical bound, at
the cost of computing at a wider intermediate width
$W \geq 2 \cdot \lceil\log_2 D\rceil / 64$.  The result is
identical to the classical HNF in~$\mathbb{Z}$.

\paragraph{The modulus must be a lattice member, not just a
covolume multiple.}
The Cohen 2.4.8 construction implicitly takes
$L \cup D\cdot\mathbb{Z}^n$ and returns its HNF, which equals
HNF$(L)$ exactly when $D \cdot \mathbb{Z}^n \subseteq L$.  In
particular, for left ideals~$I \subseteq \mathcal{O}_0$, the
reduced norm $N(I)$ has $N(I) \cdot \mathcal{O}_0 \subseteq I$
and is therefore admissible; arbitrary multiples of $\det(L)$
need not satisfy $D \cdot e_i \in L$ and produce a strictly
coarser super-lattice.

An implementor who reads ``$D$ is a multiple of $\det(L)$''
without the membership condition can derive a closed-form
covolume formula like $D = c \cdot d^4 \cdot N^2 \cdot p$ for
some scalar~$c$, hand that to the modular HNF, and silently
construct $L + D\cdot\mathbb{Z}^n$ instead of~$L$.  In the
\textsc{SuitableIdeals} cross-order step
($I^{(t)} = \mathrm{conj}(I_{\mathrm{reduced}}) \cdot J_t$), the
super-lattice may contain integer scalars like $(2, 0, 0, 0)$
that should not be in $I^{(t)}$.  Those scalars then enter the
L\textsuperscript{2}-reduced class Gram with a zero
diagonal, and the Lov\'asz/swap loop fails to terminate --- on
both our implementation and the C reference's
\texttt{quat\_lll\_core}, which we confirmed by extracting the
degenerate Gram and feeding it through C-ref directly.

The C reference avoids this by computing
$D := |\det(\text{first 4 product generators})|$ at runtime
(\texttt{lattice.c:231--233}), which is the covolume of those
generators and therefore lies in the lattice they span.
\texttt{Lattice::product} (without an explicit modulus)
implements the same recipe in Rust.  A precomputed covolume
bound is correct only when the implementer can prove
$D \cdot e_i \in L$ for the input shape --- a property that
covolume formulas do not capture.

For NIST-I commitment ideals, $D = 4 d^4 N^2 p \approx 2^{1282}$
is a compile-time constant.  For response-phase and post-reduce
ideals the modulus varies and is computed per call.

\paragraph{Recommendation.}
Add a remark after \algo{3.2} noting that fixed-precision
implementations must account for coefficient growth, and that the
standard modular HNF technique works only when the modulus~$D$
satisfies $D \cdot \mathbb{Z}^n \subseteq L$ --- not merely that
$D$ divides~$\det(L)$.  State the safe defaults explicitly:
either $D := N(I)$ (always admissible for an ideal of norm $N(I)$)
or $D := |\det(\text{first }n\text{ generators})|$ (the C
reference's choice for lattice products).

% =====================================================================
\section{Endomorphism image $P{-}Q$ and the Okeya-Sakurai lift}
\label{sec:endo-pmq}

\finding{Affects \algo{3.15}
(\textsc{FixedDegreeIsogeny}), and transitively
\textsc{IdealToIsogeny} and \textsc{KeyGen}.}{\high}

\algo{3.15} line~4 says to construct the kernel from
$(\theta(P),\, \theta(Q))$.  The $(2,2)$-isogeny chain lifts
these Montgomery $x$-only points to Jacobian $(x,y,z)$
coordinates via the Okeya-Sakurai formula, which requires the
projective $x$-only representative of $P{-}Q$ (not just its
affine $x$-coordinate).

The Okeya-Sakurai formula recovers $y_Q$ from
$(x_P, y_P, X_Q, Z_Q, X_{P{-}Q}, Z_{P{-}Q})$.  The specific
projective representative $(X_{P{-}Q} : Z_{P{-}Q})$ determines
the recovered $y_Q$.  Two representatives that agree in the
ratio $X/Z$ but differ in the individual coordinates produce
different $y_Q$ values---one correct, one not.

The spec does not specify how $\theta(P{-}Q)$ is obtained.
Two natural approaches:
\begin{enumerate}[nosep]
  \item Compute $\theta(P{-}Q)$ as a third biladder call with
        scalars $(m_{00}{-}m_{01},\, m_{10}{-}m_{11})$ on the
        same basis triple $(P, Q, P{-}Q)$.
  \item Compute \textsc{DifferencePoint}$(\theta(P),\, \theta(Q))$,
        which takes an $\Fpii$ square root and is ambiguous
        ($P{-}Q$ vs.\ $P{+}Q$).
\end{enumerate}

Approach~(1) produces three outputs whose projective
representatives are consistent (all derived from the same
basis via differential addition), so the Okeya-Sakurai
lift recovers the correct $y$.  Approach~(2) picks a
deterministic but potentially wrong branch of the square root,
causing the lift to recover the wrong $y$ for roughly half of
all inputs.  In our implementation, approach~(2) caused the
$(2,2)$-chain to fail for every $\theta$ from
\textsc{RepresentInteger} while succeeding for the $i$
automorphism (a degenerate case whose square root happens to
land on the correct branch).

\paragraph{Recommendation.}
Add a remark in \algo{3.15} or in the description of the
$(2,2)$-isogeny chain that the difference $\theta(P){-}\theta(Q)$
must be computed by applying $\theta$ to $P{-}Q$ directly (via
the biladder with the same basis triple), \emph{not} by
recomputing \textsc{DifferencePoint} on the outputs.  The same
constraint applies to every other site that constructs a basis
triple for the Okeya-Sakurai lift.

% =====================================================================
\section{\textsc{GetIndexSplitting} must handle malformed input}
\label{sec:splitting-panic}

\finding{Affects \algo{8.42}
(\textsc{GetIndexSplitting}), and transitively verification
(\algo{4.9}).}{\medium}

\algo{8.42} searches the null point $U_{i,j}(0)$ for the unique
index $(i,j)$ where it vanishes.  The spec states this index
``exists and is unique'' (Proposition~8.5.3), which is true for
honestly generated signatures.

\paragraph{What goes wrong.}
For \emph{malformed} signatures (corrupted \texttt{E\_aux},
\texttt{M\_chl}, or \texttt{chl}), the $(2,2)$-isogeny chain may
produce a null point with zero or multiple vanishing entries.
An implementation that asserts uniqueness (e.g., via
\texttt{debug\_assert!}) will panic on such inputs, providing a
denial-of-service vector: any untrusted signature can crash the
verifier in debug builds.  In release builds, the assertion compiles
out and the function returns a wrong index, leading to silent
misbehavior rather than a clean rejection.

We discovered this via negative test vectors
(\S\ref{sec:splitting-panic}): single-bit flips in
\texttt{E\_aux}, \texttt{M\_chl}, and \texttt{chl} all triggered
the assertion.

\paragraph{Recommendation.}
Note in \algo{8.42} that implementations must handle the case where
no (or multiple) vanishing indices exist, returning an error rather
than assuming uniqueness.  A verification implementation should treat
a missing or ambiguous splitting index as signature rejection, not an
assertion failure.

%% ---------------------------------------------------------------
\subsection*{Signing key encoding requires caching the ideal generator}
\label{sec:sk-generator}

\paragraph{Severity.} Low (implementation note).

\paragraph{Affected algorithm.} Key encoding (\S4.6).

\paragraph{Current spec text.}
The signing key wire format encodes the generator~$\alpha$ of the
secret ideal $I_{\mathrm{sk}} = O_0\langle\alpha, N\rangle$ as four
signed 32-byte coordinates.  The ideal is defined by~$\alpha$ and
the norm~$N$; $\alpha$ is part of the key.

\paragraph{What goes wrong.}
An implementation that converts the ideal to HNF for internal use
(as we do for lattice operations) loses the original generator.
Recovering~$\alpha$ from HNF requires a brute-force search over
small linear combinations of basis vectors.  The spec does not
mention this recovery problem.

\paragraph{Recommendation.}
Note in \S4.6 that implementations should cache~$\alpha$ at
parse/keygen time rather than attempting to recover it from the
lattice representation.

% =====================================================================
\section{\textsc{IdealToIsogeny} line 6: matrix formula is under-specified}
\label{sec:id2iso-line6}

\finding{Affects \algo{3.13} (\textsc{IdealToIsogeny}), line~6.}{\high}

\algo{3.13} line~6 writes
\[
  [P, Q]^T \leftarrow \tfrac{1}{\mathrm{nrd}(I)\,\mathrm{nrd}(J_t)}\,
    M_{\beta_1^{-1}}\, M_{\beta_2}\, [\varphi_v(P_t), \varphi_v(Q_t)]^T.
\]
The $\beta_1^{-1}$ factor is a quaternion inverse, not a matrix
inverse.  The identity derived in the surrounding text,
\[
  \beta_2 \circ \varphi_{J_t} \circ \widehat{\varphi_{J_s}} \circ
    \widehat{\beta_1}
  = [\mathrm{nrd}(I)\,\mathrm{nrd}(J_s)\,\mathrm{nrd}(J_t)]
    \widehat{\varphi_{I_2}} \circ \varphi_{I_1},
\]
uses the conjugate~$\widehat{\beta_1}$, and $\beta_1^{-1} =
\widehat{\beta_1}/\mathrm{nrd}(\beta_1)$ only as quaternions.  An
implementor writing $M_{\beta_1^{-1}}$ as a \emph{matrix inverse}
gets a different matrix from the intended
$M_{\widehat{\beta_1}}/\mathrm{nrd}(\beta_1) =
\mathrm{adj}(M_{\beta_1})/\mathrm{nrd}(\beta_1) \pmod{2^f}$ in
general.

\paragraph{What the C reference does.}
\texttt{dim2id2iso.c:881--1014} builds the quaternion
$\theta = \beta_2 \cdot \widehat{\beta_1}$ directly and scales its
coordinates by $\mathrm{invmod}(d_1 \cdot \mathrm{nrd}(J_t),
2^f)$ before decomposing on~$\mathcal{O}_0$.  This is
$\beta_2 \cdot \widehat{\beta_1}$, whereas the spec's
$\beta_1^{-1}\,\beta_2 = \widehat{\beta_1}\,\beta_2 /
\mathrm{nrd}(\beta_1)$ has the factors in the opposite order; the
two quaternion products agree only after the subsequent
$\mathrm{nrd}(\beta_1)$ scaling, and quaternion multiplication is
not commutative in general.

\paragraph{Impact.}
An implementor has two independent traps: (a) interpreting
$M_{\beta_1^{-1}}$ as the matrix inverse rather than the action
matrix of~$\beta_1^{-1}$ viewed as a quaternion; (b) writing the
quaternion product in the opposite order from the derivation.

\paragraph{Recommendation.}
Write line~6 with the conjugate factored out and the scaling
collected:
\[
  [P, Q]^T \leftarrow
    \tfrac{1}{\mathrm{nrd}(I)\,\mathrm{nrd}(J_t)\,\mathrm{nrd}(\beta_1)}
    M_{\widehat{\beta_1}}\, M_{\beta_2}\, [\varphi_v(P_t),
    \varphi_v(Q_t)]^T,
\]
and state which quaternion-product order the matrix composition
uses (equivalently, specify whether $M_{\alpha\beta} = M_\alpha
M_\beta$ or $M_\beta M_\alpha$).

% =====================================================================
\section{\textsc{IdealToIsogeny}:
  \texorpdfstring{$\beta$}{beta} lives in the reduced ideal, not the
  input \texorpdfstring{$I$}{I}}
\label{sec:id2iso-reduced-ideal}

\finding{Affects \algo{3.13}, \algo{3.16} (\textsc{SuitableIdeals}),
  and anywhere the scaling $1/\mathrm{nrd}(I)$ appears downstream
  of short-vector enumeration.}{\high}

\algo{3.16} (\textsc{SuitableIdeals}) returns $\beta_i$ obtained
from the short-vector enumeration inside $J_i\,I$, and~$d_i =
\mathrm{nrd}(\beta_i)/\mathrm{nrd}(J_i\,I)$.  Efficient
implementations (and the C reference,
\texttt{dim2id2iso.c:534--565}) first replace~$I$ by the
\emph{smallest equivalent ideal}~$I' = I \cdot
\widehat{\delta}/\mathrm{nrd}(I)$, so that the short-vector
search produces usable degrees.  The returned $\beta$ and~$d$ are
relative to~$I'$, not to~$I$, so the downstream invariant is
\[
  \mathrm{nrd}(\beta) = d \cdot \mathrm{nrd}(I')
  \neq d \cdot \mathrm{nrd}(I).
\]

\paragraph{What the C reference does.}
\texttt{dim2id2iso.c:670--684} ``sends~$\beta$ back'' to the
original ideal after the pair is selected (multiplying
by~$\delta$), so that by the time \algo{3.13} line~6 applies the
scaling $1/\mathrm{nrd}(I)$, the identity
$\mathrm{nrd}(\beta_1) = d_1 \cdot \mathrm{nrd}(I)$ holds in
terms of the caller-supplied~$I$.

\paragraph{What goes wrong.}
An implementor who keeps~$\beta$ in the reduced ideal~$I'$ (a
natural data representation) must also use $\mathrm{nrd}(I')$ in
every downstream scaling formula.  Using the spec's
$\mathrm{nrd}(I)$ literal gives $\det(M_{\beta_1}) \neq d_1 \cdot
\mathrm{nrd}(I) \pmod{2^f}$, silently producing a non-isotropic
outer kernel and \texttt{zeros=0} in the final splitting step.

\paragraph{Recommendation.}
State in \algo{3.16} that the returned~$\beta$ must be expressed
in the input ideal's lattice (so $\mathrm{nrd}(\beta) = d \cdot
\mathrm{nrd}(I)$), or name the lattice that~$\beta$ lives in and
refer to its norm throughout \algo{3.13}.

% =====================================================================
\section{\textsc{SuitableIdeals} line 14:
  \texorpdfstring{$v_2(u)$}{v2(u)} vs.\
  \texorpdfstring{$v_2(\gcd(u, v))$}{v2(gcd(u,v))}}
\label{sec:suitableideals-e-val}

\finding{Affects \algo{3.16} (\textsc{SuitableIdeals}), line~14.}{\medium}

\algo{3.16} line~14 writes ``$e \leftarrow v_2(u)$'' and line~15
divides both~$u$ and~$v$ by~$2^e$.  Using $v_2(u)$ is only
equivalent to $v_2(\gcd(u, v))$ when $v_2(u) \le v_2(v)$:
\begin{itemize}
  \item If $v_2(u) > v_2(v)$, dividing~$v$ by~$2^{v_2(u)}$ loses
    low bits of~$v$ that were carrying real value.
  \item If $v_2(u) < v_2(v)$, dividing by~$2^{v_2(u)}$ leaves~$v$
    still even, so the next algorithm's precondition
    $\gcd(u \cdot d_1, v \cdot d_2) = 1$ fails.
\end{itemize}

\paragraph{What the C reference does.}
\texttt{dim2id2iso.c:833} computes
\texttt{ibz\_two\_adic(\&gcd(u, v))}, i.e.\ $v_2(\gcd(u, v))$.

\paragraph{Recommendation.}
Write line~14 as ``$e \leftarrow v_2(\gcd(u, v))$''.

% =====================================================================
\section{\algo{8.47} implementation variants are not specified}
\label{sec:iso22chain-variants}

\finding{Affects \algo{8.47} (\textsc{Isogeny22ChainWithTorsion})
  and its callers, primarily \algo{3.13} line~9.  Observed
  $99\%$-wrong-codomain failure rate in keygen before the
  implementation accounted for the kernel-torsion convention.}{\high}

The spec describes a single uniform chain of $e$ 8-torsion
isogenies, with the phrasing ``a chain of $e$ $(2,2)$-isogenies of
2-power degree.''  It does not make explicit either
\begin{enumerate}
  \item what kernel torsion the chain consumes, nor
  \item that \algo{3.13} and \algo{3.15} call the chain with
    different conventions.
\end{enumerate}
The C reference implementation in
\texttt{theta\_isogenies.c} splits \algo{8.47} into two variants
keyed on a boolean \texttt{extra\_torsion}:
\begin{itemize}
  \item \texttt{extra\_torsion = true}: an $e$-length chain of
    8-torsion isogenies.  Kernel must have order~$2^{e+2}$.  The
    last two chain steps use the \texttt{hadamard\_bool}
    convention from \algo{8.41} to fold the kernel's 4-torsion
    and 2-torsion residues into the penultimate and ultimate
    $(2,2)$-isogenies.
  \item \texttt{extra\_torsion = false}: an $(e-2)$-length chain
    of 8-torsion isogenies, followed by a dedicated 4-isogeny
    step (\texttt{theta\_isogeny\_compute\_4}) and a dedicated
    2-isogeny step (\texttt{theta\_isogeny\_compute\_2}).
    Consumes a kernel of order~$2^e$.  The final two steps
    compute square roots rather than relying on extra torsion.
\end{itemize}
\algo{3.13} line~9 (the outer chain) calls the implementation
with \texttt{extra\_torsion = false}
(\texttt{dim2id2iso.c:1128}); the inner \algo{3.15} calls it with
\texttt{extra\_torsion = true} (\texttt{dim2id2iso.c:181}).  The
two variants produce the same abstract~$(2^e, 2^e)$-isogeny but
require different kernel torsions.

\paragraph{Concrete impact (Selkie experience).}
We implemented the 8-torsion chain corresponding to
\texttt{extra\_torsion = true} and used it for every call.
\algo{3.13} line~9 passes a kernel whose natural order is
$2^e$ (from doubling the $2^f$-torsion image basis by $f-e$),
so our chain received a kernel that was two torsion bits short of
what the \texttt{extra\_torsion = true} path expects.  The chain
ran to completion in every instance, the splitting code
(\algo{8.41}) silently selected one of the ten $U_{i,j}(0)$
indices, and keygen returned a curve that is not the public
key~$E_\mathrm{pk}$.  $100\%$ of outer chains produced
terminal splits with either~$0$ or~$10$ vanishing $U_{i,j}(0)$
coordinates rather than the unique~$1$ that identifies a genuine
elliptic product.  Downstream code paths did not detect the
failure: \algo{4.3}, \algo{4.6}, and \algo{4.9} all accept
whichever curve the chain produces.

We resolved this by over-padding the kernel: instead of
doubling~$f-e$ times, we double~$f-e-2$ times, so the kernel
enters the chain with order~$2^{e+2}$ and the
\texttt{extra\_torsion = true} path works unchanged.  This
requires $e \le f-2$.  The side condition reduces to
$v_2(\gcd(u,v)) \ge 2$ in \textsc{find\_uv\_from\_lists}, which
we filter for in \algo{3.16}.

\paragraph{Why this is hard to catch in testing.}
\begin{itemize}
  \item The chain does not fault when given a short kernel; it
    returns a malformed abelian surface.
  \item The splitting routine \algo{8.42}
    (\textsc{SplittingIsomorphism}) picks a codomain side by
    locating a zero among the ten~$U_{i,j}(0)$ coordinates and
    does not check that exactly one vanishes.  When all ten
    vanish (the fully degenerate case), any of the ten branches
    is ``valid''; the spec doesn't say which to pick.  When none
    vanish, the function should error but the text only says
    ``find the unique index such that~$U_{i,j}(0) = 0$.''
  \item Downstream consumers of the split codomain (basis
    computation in \algo{2.3}, Tate pairing in \algo{2.5},
    kernel-to-ideal conversion in \algo{3.17}) accept
    non-product curves without visible failure.  They compute
    \emph{something}, but it's not the object the protocol
    calls for.
  \item Only byte-for-byte comparison with KAT vectors reveals
    the divergence; unit tests on synthetic kernels with small
    $e$ do not exhibit the bug because $f - e \gg 2$ leaves
    enough spare torsion implicitly.
\end{itemize}

\paragraph{Recommendation.}
We suggest three spec-level changes, in decreasing order of
priority:
\begin{enumerate}
  \item \algo{8.47} should state the kernel-torsion requirement
    explicitly.  ``A chain of~$e$ $(2,2)$-isogenies consumes a
    kernel of order~$2^{e+2}$'' (current implicit convention) or
    ``of order~$2^e$'' (shorter variant)
    unambiguously pins down the caller's contract.
  \item \algo{8.41} (\textsc{SplittingIsomorphism}) or
    \algo{8.42} should add a precondition:
    \emph{exactly one} $U_{i,j}(0)$ must vanish at the chain's
    terminal theta null.  An implementer who checks this
    invariant gets a clear failure signal instead of silent
    corruption.
  \item \algo{3.13} line~9 should state explicitly which chain
    variant it uses, so that an implementer who supplies only
    one variant knows the coupling fails and either adapts
    (pad kernel up, pad chain down) or implements the
    missing variant.
\end{enumerate}
Absent these changes, a faithful implementation that reads the
spec straight through will silently produce invalid signing keys
that fail only when compared against known-answer test vectors.

% =====================================================================
\section{\textsc{IdealToIsogeny} output basis must be aligned with
  the isogeny, not canonical on the codomain}
\label{sec:id2iso-aligned-basis}

\finding{Affects \algo{3.13} (\textsc{IdealToIsogeny}) and
  \algo{4.5} (\textsc{SplitAuxiliaryIsogeny}).}{\high}

\algo{3.13} states that its output is ``the codomain~$E_I$
of~$\varphi_I$, $\varphi_I(P_0)$, and $\varphi_I(Q_0)$.'' Read
literally, this is the image of the canonical $\mathcal{O}_0$
basis~$(P_0, Q_0)$ of~$E_0[2^f]$ through the full
ideal-to-isogeny~$\varphi_I$.

\paragraph{The implicit requirement.} \algo{4.5}
(\textsc{SplitAuxiliaryIsogeny}) takes as input curves~$E_1$,
$E_2$ and bases~$(P_1, Q_1)$, $(P_2, Q_2)$ together with
``positive integers~$q$, $e'$, $r$ such that there exist four
isogenies $\varphi_q$, $\varphi_{2^{e'}-q}$, $\psi_{2^{e'}-q}$,
$\psi_q$ verifying \dots\ $(P_2, Q_2) = (\varphi(P_1),
\varphi(Q_1))$.''  The condition $(P_2, Q_2) = \varphi(P_1, Q_1)$
forces $(P_2, Q_2)$ to be the image of $(P_1, Q_1)$ through the
specific odd-degree isogeny $\varphi$ of degree~$q\,(2^{e'}-q)$.
An independently-chosen torsion basis on $E_2$ does not satisfy
this and produces a kernel that fails the splitting condition
inside \algo{8.47}.

\paragraph{Why this is implicit.}
The spec states the output of \algo{3.13} as ``$\varphi_I(P_0)$,
$\varphi_I(Q_0)$.'' The composition~$\varphi_I = \varphi_{2^e} \circ
\dots \circ \varphi_d$ has both odd-degree and~$2$-power factors.
The image of~$P_0$ through the full composition retains
$2^f$-torsion only because every odd-degree factor preserves the
$2^f$-torsion subgroup; the~$2$-power factors are evaluated
through the dimension-2~$(2,2)$-chain on the abelian-surface
side, where the input torsion is determined by what the
implementer pushes through the chain.  The C reference's
\textsc{Clapotis} algorithm
(\texttt{dim2id2iso\_ideal\_to\_isogeny\_clapotis} in
\texttt{dim2id2iso.c:778--1228}) navigates this by pushing the
canonical basis of the \emph{intermediate} extremal-order
curve~$E_t$ through the~$(2,2)$-chain after the
\textsc{FixedDegreeIsogeny} step has already produced an aligned
basis on the chain's input curve, $E_u$.  The final output basis
on $E_I$ is then the image of $E_t$'s basis (not~$E_0$'s) through
the chain, which is consistent with the alignment required
by~\algo{4.5}.  None of this is in the spec.

\paragraph{Impact.}
A naive implementation following \algo{3.13} returns either
(a)~the literal $\varphi_I(P_0), \varphi_I(Q_0)$ via push-through
of $E_0$'s basis through the entire chain --- the resulting points
have torsion $2^{f - 2\,\mathsf{chain\_e}}$, which is trivial for
realistic parameters where $\mathsf{chain\_e}$ approaches $f$ ---
or (b)~a freshly-sampled canonical $2^f$-torsion basis on~$E_I$
unrelated to~$\varphi_I$.  Either choice causes \algo{4.5} to
silently fail at the~$(2,2)$-chain splitting check.

\paragraph{What the C reference does.}
The C reference's \textsc{Clapotis}
(\texttt{dim2id2iso\_ideal\_to\_isogeny\_clapotis}) effectively
treats the ideal-to-isogeny output basis as ``aligned with~$\varphi_I$
in a way that downstream~\algo{4.5} can consume.''  Concretely it
takes the canonical basis of the matching extremal-order
curve~$E_t$, pushes it through \textsc{FixedDegreeIsogeny} (an
odd-degree isogeny~$\varphi_v$ that preserves $2^f$-torsion),
applies the endomorphism~$\theta = \beta_2\,\widehat{\beta_1}$ to
get a kernel structure on~$E_u \times E_v$, runs the~$(2,2)$-chain
with the basis-preserving push-through, and reports the surviving
basis on the chain's first codomain factor.

\paragraph{Recommendation.}  The spec should:
\begin{enumerate}
  \item State explicitly that the output of \algo{3.13} is a
    pair $(R, S) \in E_I[2^f]^2$ such that $(R, S) = \nu(P_0, Q_0)$
    for an isogeny~$\nu$ \emph{whose 2-adic valuation} the caller
    can recover from $(\beta_1, \beta_2, u, v, d_1, d_2)$, and
    pin down the convention so~\algo{4.5}'s precondition
    $(P_2, Q_2) = \varphi(P_1, Q_1)$ is realized.
  \item Reference the C reference's basis-alignment construction
    (or describe it inline) so an implementor doesn't read the
    \emph{output} of \algo{3.13} as ``the literal image of $E_0$'s
    basis through the full chain'' --- which has the wrong
    torsion --- nor as ``a fresh canonical basis on $E_I$'' ---
    which has the right torsion but no alignment with~$\varphi$.
\end{enumerate}
We follow the C reference's basis-alignment construction in our
implementation; the resulting basis is what \algo{4.5} consumes
without further adjustment.

% =====================================================================
\section{\textsc{IdealToIsogeny} must propagate $P{-}Q$ alongside
  $(P, Q)$}
\label{sec:id2iso-pmq}

\finding{Affects \algo{3.13} (\textsc{IdealToIsogeny}) output and
  \algo{4.5} (\textsc{SplitAuxiliaryIsogeny}) input.}{\critical}

\algo{3.13} writes its output as ``$E_I, \varphi_I(P_0),
\varphi_I(Q_0)$''---a curve and two points, no third
coordinate.  But \algo{4.5} forms a $(2, 2)$-kernel from those
two points (line~3, ``$T \leftarrow [2^{f-e'}](\varphi_{\text{com}}(P_0),
[q_{\text{rsp}}^{-1}\,2^{f-e'}]\varphi_{\text{com,rsp,aux}}(P_0))$''),
and any $(2, 2)$-chain that consumes such a kernel must call
\textsc{lift\_basis} (Okeya--Sakurai) to recover the Jacobian
$y$-coordinate.  Okeya--Sakurai requires a third basis
component~$P{-}Q$ \emph{whose projective representative is
consistent with the chain's prior evaluations of $P$ and~$Q$} ---
see~\S\ref{sec:pmq-recompute} and~\S\ref{sec:endo-pmq} for the
already-noted instances of the same pitfall.

If \algo{3.13} returns only $(\varphi_I(P_0), \varphi_I(Q_0))$
and the next stage recovers $\varphi_I(P_0 - Q_0)$ via
\textsc{DifferencePoint}, the fresh square root branches
independently of the chain's internal lift.  The downstream
$(2, 2)$-chain (here the response-phase chain inside
\algo{4.5}) then reaches its terminal theta null with
$\#\{i, j : U_{i, j}(0) = 0\} = 0$ instead of the unique~$1$
that identifies a product splitting (\S\ref{sec:get-index}),
and the chain fails.

\paragraph{What the C reference does.}
A C-reference \texttt{theta\_couple\_curve\_with\_basis\_t}
always carries~$B.\mathrm{PmQ}$ alongside $B.P$ and~$B.Q$, and
every isogeny evaluation pushes all three through together.
\texttt{dim2id2iso\_ideal\_to\_isogeny\_clapotis}
(\texttt{dim2id2iso.c:778--1228}) does so for the
\algo{3.13} chain, and the \texttt{copy\_bases\_to\_kernel}
helper inside \texttt{compute\_dim2\_isogeny\_challenge}
(\texttt{sign.c:240--256}) carries~$\mathrm{PmQ}$ into the
response-phase kernel.  No conversion to/from
$\textsc{DifferencePoint}$ ever occurs after the canonical
basis is lifted.

\paragraph{Impact.}
A literal reading of \algo{3.13}'s output triple
$(E_I, \varphi_I(P_0), \varphi_I(Q_0))$ leads an implementor to
recover $P{-}Q$ via \textsc{DifferencePoint} at the next call
site, which silently corrupts \algo{4.5} on every input, with
no error message: the chain returns ``no splitting found''
which is indistinguishable from a probabilistic chain failure
that should trigger a retry.  The implementor concludes the
input ideals are pathological rather than that the basis is
malformed.

\paragraph{Recommendation.}
Make explicit in \algo{3.13} that the output is a triple
$(R, S, R{-}S) \in E_I[2^f]^3$, not a pair, and that the third
component must be the propagated image~$\varphi_I(P_0 - Q_0)$ of
the canonical basis difference --- never recovered via
\textsc{DifferencePoint} from~$R$ and~$S$.  Add a corresponding
note to \algo{4.5} that its input bases on~$E_1$ and~$E_2$ each
include a propagated~$\mathrm{PmQ}$, and that the kernel
construction (lines~1--4) must scale~$\mathrm{PmQ}$ alongside
$P$ and~$Q$ rather than recomputing.  Equivalently, update the
\textsc{TorsionBasisFromHint} and \textsc{IdealToIsogeny} signatures
to return the standard $(P, Q, P{-}Q)$ triple that \S2 already
implicitly assumes throughout.

% =====================================================================
\section{\textsc{SplittingIsomorphism} omits a side-channel
randomization that the published KAT vectors require}
\label{sec:specrev-splitter-randomization}

\finding{Affects \algo{8.42} (\textsc{SplittingIsomorphism}) and
  the chain output of any algorithm that consumes it (\algo{3.13},
  \algo{4.5}).}{\high}

The output of an extra-torsion-free $(2, 2)$-isogeny chain is a
level-2 theta null point on a product abelian surface
$E_1 \times E_2$.  Multiple projective representatives exist for
the same abstract surface, related by the symplectic
modular-group action.  \algo{8.42} as written picks one such
representative deterministically from the chain's terminal theta
null.  But the choice is a function of the secret kernel, so the
emitted projective representative leaks information about the
kernel beyond what the curve isomorphism class already exposes.

\paragraph{What the C reference does.}
The C reference closes this leak in
\texttt{splitting\_compute} (\texttt{theta\_isogenies.c:1001--1069})
by pre-multiplying the deterministic splitting matrix~$M$ by a
uniformly random $M_k$ drawn from a precomputed list of six
$4 \times 4$ matrices over~$\Fpii$.  The six matrices are tensor
products $M_k = m_k \otimes m_k$ of $2 \times 2$ coset
representatives of the symplectic-modular-group quotient
$\Gamma / \Gamma^0(4)$ that, as a group, suffice to mix the
representative within its equivalence class when working in the
Montgomery model.  The matrices are built by
\texttt{precompute\_hd\_splitting.sage} and shipped as
\texttt{NORMALIZATION\_TRANSFORMS} in
\texttt{precomp/.../hd\_splitting\_transforms.c}.

The index~$k \in \{0, \dots, 5\}$ is sampled by reading four bytes
from the deterministic byte-replayable RNG, parsing them as a
little-endian \texttt{u32}, rejection-sampling against threshold
$4{,}294{,}967{,}292 = 6 \cdot \lfloor 2^{32} / 6 \rfloor$ for an
unbiased~$\bmod\,6$, and indexing the precomputed list
(\texttt{theta\_isogenies.c:980}, \texttt{sample\_random\_index}).
The randomization is gated on \texttt{\#if defined(ENABLE\_SIGN)}
and on a caller-supplied \texttt{randomize} flag: it is on for the
keygen and signing outer chains
(\texttt{theta\_chain\_compute\_and\_eval\_randomized}), and off
for verification (\texttt{theta\_chain\_compute\_and\_eval\_verify})
and for the deterministic \textsc{FixedDegreeIsogeny} chains
(\texttt{theta\_chain\_compute\_and\_eval}).

\paragraph{Two consequences.}

\begin{enumerate}

\item It is a real side-channel hardening.  Without it, the public
  output of keygen and signing is a deterministic function of the
  secret kernel structure, not just the abstract codomain.  An
  attacker observing the projective representative of a public
  key, signature, or commitment curve could potentially recover
  information about the underlying kernel beyond what the curve
  isomorphism class already reveals.  Whether the leakage is
  exploitable in practice for SQIsign is, to our knowledge, an
  open question, but the cost of the countermeasure is six
  precomputed matrices, four bytes of RNG, and one matrix
  multiplication.

\item The published KAT vectors are produced \emph{with} this
  randomization applied.  An implementation that follows the spec
  literally and skips the randomization will produce a public key
  that is projectively equivalent to the KAT vector's, but
  byte-different.  Without mirroring the exact same RNG byte
  consumption pattern (4 bytes, little-endian, threshold
  $4{,}294{,}967{,}292$, $\bmod\,6$) and matrix selection logic,
  byte-equality with the published KATs is impossible.  We spent
  significant effort tracing chain-internal byte mismatches that
  vanished only after replicating the splitter randomization.

\end{enumerate}

\paragraph{Recommendation.}
Add an explicit \textsc{SplittingNormalization} step to
\algo{8.42} (or as a post-processing step on its output) that
specifies:

\begin{itemize}
\item The list of six normalization matrices $M_0, \dots, M_5$
  (or a generating procedure --- the Sage script that produces
  them is short).
\item The exact RNG byte consumption pattern (four bytes,
  little-endian \texttt{u32}, rejection threshold
  $4{,}294{,}967{,}292$, $\bmod\,6$).
\item A precise statement of which chain calls (keygen and
  signing outer chains, but not FDI sub-chains and not
  verification) randomize, and which use the deterministic form.
\item The threat model: that the deterministic representative
  leaks kernel structure, and that the randomization is required
  to prevent that leak.
\end{itemize}

Without these specifics, an implementation that follows the spec
correctly will still fail the KAT comparison and the implementor
will have no way to discover the discrepancy without reading the
C reference's \texttt{splitting\_compute}.

% =====================================================================
\section{\textsc{SuitableIdeals} multi-order optimization: the
  cross-order $\beta$ transport step}
\label{sec:specrev-suitable-multi-order-transport}

\finding{Affects \algo{3.16} (\textsc{SuitableIdeals}) when extended
  with the multi-order optimization, and any caller that expects
  $\beta_i \in I$ (the input ideal).  Causes \algo{3.13}
  (\textsc{IdealToIsogeny}) to silently fail
  when the special-order pair is unavailable.}{\high}

\algo{3.16} as written enumerates short vectors only in the input
ideal~$I$ (or its LLL-reduced equivalent~$I'$).  When this
single-lattice search exhausts without finding a viable
$(\beta_s, \beta_t)$ pair, the C reference falls back to the
\emph{multi-order} optimization
(\texttt{dim2id2iso.c:534--666}): it enumerates short vectors in
six additional pushforward lattices
\[
  L_t \;=\; \overline{I'} \cdot J_t,
  \qquad t = 1, \ldots, 6,
\]
where $J_t$ is the precomputed connecting ideal from $\mathcal{O}_0$
to a maximal order $\mathcal{O}_t$ (the right order of $J_t$ is
conjugate to~$\mathcal{O}_t$).  This more-than-doubles the
candidate pool and rescues KATs whose hypercube enumeration would
otherwise exhaust.

A short vector $\beta_t \in L_t$ is an element of
$\overline{I'} \cdot J_t \subseteq B$, \emph{not} of~$I$.  Every
downstream caller of \textsc{SuitableIdeals} (\algo{3.13},
\algo{3.17}, \algo{4.6}) expects the returned $\beta_i$ to live
in the original ideal lattice~$I$ so they can compute
$\theta = \beta_2 \cdot \overline{\beta_1} / \mathrm{nrd}(I)$ and
treat it as an endomorphism of the protocol's reference curve.
The multi-order optimization therefore needs an explicit transport
step from $L_t$ back to~$I$ before returning $\beta_i$.

\paragraph{What the C reference does.}
\texttt{dim2id2iso.c:651--672}, after \texttt{find\_uv\_from\_lists}
selects an $(i_1, i_2)$ pair, applies the following transformation
when either index $j_1$ or $j_2$ is non-zero (i.e., the chosen
$\beta$ came from an alternate-order lattice):
\begin{enumerate}
  \item Adjust \texttt{delta.denom}, where \texttt{delta} is the
    conjugated LLL-first vector $\overline{\delta}$ of $I'$:
\begin{verbatim}
ibz_div(&delta.denom, &remain, &delta.denom, &lideal->norm);
ibz_mul(&delta.denom, &delta.denom, &conj_ideal.norm);
\end{verbatim}
    so that \texttt{delta} ends up at denom $d_{I'} \cdot
    \mathrm{nrd}(\overline{I'})$.
  \item For each $j_i \neq 0$:
\begin{verbatim}
quat_alg_mul(beta_i, &delta, beta_i, Bpoo);
quat_alg_normalize(beta_i);
quat_alg_conj(beta_i, beta_i);
\end{verbatim}
\end{enumerate}
Algebraically: $\beta_i \leftarrow \overline{\,\overline{\delta}
\cdot \beta_i\,}$.  The post-condition asserted at
\texttt{dim2id2iso.c:690} is
$\beta_i \in I$ (i.e., \texttt{quat\_lattice\_contains(\&I,
beta\_i)}).

This step is not in \algo{3.16}'s pseudocode at all.  An
implementor reading the spec will write \textsc{SuitableIdeals}
without this transport, observe that 99\% of NIST-I KATs pass
(because \textsc{SuitableIdeals} usually finds a special-order
pair on the first try), and only discover the bug on a KAT that
forces the multi-order branch.  In the published NIST-I lvl1 KAT
file, that's count~29: every other count from~0 through~99 finds a
special-order pair within the
\texttt{FINDUV\_box\_size = 2} hypercube.

\paragraph{Concrete impact (Selkie experience).}
Without the transport, our $\beta_2$ from $L_1 = \overline{I'}
\cdot J_1$ flowed into the downstream
$\texttt{action\_matrix}(\beta_2, \mathcal{O}_1)$ call, whose
\texttt{decompose} step rejects $\beta_2$ as ``not in the
lattice'' and aborts the entire \algo{3.13} flow with no error
signal --- key generation transparently retries from the next
random commitment, eventually picks an $(s, t) = (0, 0)$ pair on
some later iter, and produces a public key that fails the KAT
byte comparison.

\paragraph{Recommendation.}
\algo{3.16} should be extended to describe the multi-order
optimization explicitly, including the post-search transport.
Concretely:
\begin{enumerate}
  \item \textbf{Search step.} Define
    $L_0 = I'$ and
    $L_t = \overline{I'} \cdot J_t$ for $t = 1, \ldots, 6$, where
    $J_t$ is the connecting ideal from $\mathcal{O}_0$ to
    $\mathcal{O}_t$.
    Enumerate short vectors in each $L_t$ separately; sort each
    list by reduced norm.
  \item \textbf{Pair selection.} Run the line-walk in
    \texttt{find\_uv\_from\_lists} over the cross product of
    sorted lists.  First valid $(\beta_s \in L_s,
    \beta_t \in L_t)$ wins.
  \item \textbf{Transport step.} If $s \neq 0$, replace
    $\beta_s \leftarrow \overline{\overline{\delta} \cdot \beta_s}$;
    if $t \neq 0$, replace
    $\beta_t \leftarrow \overline{\overline{\delta} \cdot \beta_t}$.
    Here $\overline{\delta}$ is the conjugate of the LLL-first
    vector of~$I'$, with denominator $d_{I'} \cdot N(\overline{I'})$
    after the modular adjustment described above.
\end{enumerate}
The transport's algebraic role is to map an
$\overline{I'} \cdot J_t$-element back to an $I$-element via the
$\overline{I'} \cdot \overline{\delta} = N(\overline{I'}) \cdot
\overline{I}$ identity (and then conjugate to flip back to a
left-$\mathcal{O}_0$ element).

The spec's silence on this transport is part of a broader pattern:
the multi-order optimization is invoked extensively by the C
reference but is not modeled in \algo{3.16} at all.  Implementations
that follow the spec literally end up with a single-lattice
\textsc{SuitableIdeals} that fails on KATs where the multi-order
search is needed.

% =====================================================================
\section{Connecting-ideal norm encoding: reduced norm vs.\ first
  basis entry}
\label{sec:specrev-connecting-ideal-norm}

\finding{Affects the precomputation of the seven extremal-order
connecting ideals~$J_t$ used by the multi-order
\textsc{SuitableIdeals} optimization (\algo{3.16}) and by
\algo{3.13} (\textsc{IdealToIsogeny}) line~6's outer-chain
scaling.}{\high}

The seven extremal maximal orders $\mathcal{O}_t$ come with
precomputed connecting ideals
$J_t \subseteq \mathcal{O}_0$ whose right orders are conjugate
to~$\mathcal{O}_t$.  Each~$J_t$ is encoded as a basis matrix
$B_t$, a denominator $d_t$, and a labelled norm $N_t$.  The
rational basis of $J_t$ is given by the columns of $B_t / d_t$,
and the labelled norm is supposed to be the
\emph{reduced norm} $N(J_t) = \sqrt{[\mathcal{O}_0 : J_t]}$.

\paragraph{What goes wrong.}
For the HNF shape used by the NIST-I precomputed ideals
($B_t = \mathrm{diag}(2 N_t, 2 N_t, 1, 1)$ at $d_t = 2$), the
reduced norm is~$N_t$ but the first-coordinate entry of the
integer basis is~$2 N_t = d_t \cdot N_t$.  An implementer who
extracts the labelled norm by reading the leading entry of the
basis ends up encoding $d_t \cdot N_t$ instead of~$N_t$,
i.e.\ the labelled norm is exactly~$d_t$ times too large for
every cross-order $t > 0$.

\paragraph{Downstream impact.}
\begin{itemize}
  \item \textsc{SuitableIdeals}' integrality check
    ($\mathrm{nrd}(\beta) / (N_t \cdot d_t^2)$) silently rejects
    half the hypercube candidates and silently halves the
    degrees of the rest.  The resulting sorted batch differs
    from the C reference's, and
    \textsc{find\_uv\_from\_lists} picks a different
    $(\beta_s, \beta_t)$.
  \item \algo{3.13} line~6's outer-chain scaling
    $1/(\mathrm{nrd}(I) \cdot d_1 \cdot N(J_t))$ for the
    cross-order case is multiplied by~$d_t$, embedding the
    output basis in a sublattice of half the intended index.
\end{itemize}
This is latent on most KATs because the special-order path
($t = 0$, which doesn't read the labelled norm at all) usually
succeeds.  KATs that force the cross-order fallback expose
both divergences in the same iteration.

\paragraph{Recommendation.}
The spec should define the connecting-ideal precomputation
format explicitly: it is a triple $(B_t, d_t, N_t)$ where
\begin{itemize}
  \item $B_t / d_t$ are the rational basis columns,
  \item $N_t$ is the reduced norm $N(J_t)$ (\emph{not} the
    leading basis entry, which equals $d_t \cdot N_t$ for the
    typical HNF shape).
\end{itemize}
The natural pitfall is to extract~$N_t$ as the first
coordinate of the integer HNF row space, which yields a
$d_t$-times-too-large value.  Implementations should derive
$N_t$ via $\sqrt{|\det(B_t)| / d_t^4}$ or equivalently
$\sqrt{[\mathcal{O}_0 : J_t]}$, matching what
\texttt{quat\_lideal\_norm} computes
(\texttt{quaternion/ref/generic/ideal.c:7}).

% =====================================================================
\section{\textsc{IdealToIsogeny} action-matrix order index after
  cross-order $\beta$ transport}
\label{sec:specrev-action-matrix-index}

\finding{Affects \algo{3.13} (\textsc{IdealToIsogeny}) when its
  caller-provided $(\beta_1, \beta_2)$ came from the multi-order
  optimization in \textsc{SuitableIdeals}.}{\medium}

After the transport described in
\S\ref{sec:specrev-suitable-multi-order-transport}, both $\beta_1$
and $\beta_2$ live in $I \subseteq \mathcal{O}_0$, regardless of
which alternate-order lattice they originated from.  An
implementor reading \algo{3.13} naturally writes
\textsc{ActionByTranslation} (line~7) to decompose $\beta_i$ on
the basis of $\mathcal{O}_{s}$ (resp.\ $\mathcal{O}_{t}$), since
the algorithm associates $\beta_1$ with the curve~$E_s$ and
$\beta_2$ with~$E_t$.  This decomposition fails: $\beta_i$ is
in~$I \subseteq \mathcal{O}_0$, not in $\mathcal{O}_s$
(or $\mathcal{O}_t$) when $s \neq 0$ (or $t \neq 0$).

\paragraph{What the C reference does.}
\texttt{endomorphism\_application\_even\_basis}
(\texttt{id2iso.c:140}) takes an explicit
\texttt{index\_alternate\_curve} argument selecting which
order's basis to decompose on.  In every call from
\textsc{IdealToIsogeny}, regardless of which alternate order
$\beta_i$ originated from, the C reference passes \texttt{0}
(see \texttt{dim2id2iso.c:1054} and \texttt{:1156}):
\begin{verbatim}
endomorphism_application_even_basis(
    &bas2, 0, &Fv_codomain.E1, &theta, TORSION_EVEN_POWER);
\end{verbatim}
i.e., always decompose on $\mathcal{O}_0$.  This is consistent
with the post-transport invariant
$\beta_i \in I \subseteq \mathcal{O}_0$.

The two conventions
(decompose on $\mathcal{O}_s$ vs.\ decompose on $\mathcal{O}_0$)
are mathematically equivalent for $\beta \in
\mathcal{O}_0 \cap \mathcal{O}_s$.  But for cross-order
$\beta$, only the $\mathcal{O}_0$ decomposition succeeds; an
implementer who reads ``$M_\beta$ is the action of $\beta$ on
$E_s$'' as ``decompose $\beta$ on $\mathcal{O}_s$'' has the
right idea for the $s = 0$ case and the wrong idea for $s \neq 0$.

\paragraph{Recommendation.}
\algo{3.13} should specify that the action-matrix decomposition
is taken on $\mathcal{O}_0$ regardless of the source order.
Equivalently: the \textsc{ActionByTranslation} subroutine should
take both the curve and the order as arguments, and \algo{3.13}
should always pass $(\mathcal{O}_0, E_*)$ for some target
curve~$E_*$.  A one-line addition of the form
\begin{quote}\itshape
  In line~7, the action-matrix decomposition
  $M_{\beta_i} = \sum_k c_k\, M_{b_k}$ is taken on the basis
  $(b_0, \ldots, b_3)$ of $\mathcal{O}_0$, where $\beta_i =
  \sum_k c_k b_k$ (which is well-defined because the
  multi-order \textsc{SuitableIdeals} returns $\beta_i \in I
  \subseteq \mathcal{O}_0$).
\end{quote}
suffices.

\paragraph{Why this is hard to catch.}
The two-decomposition factoring
$M_{\beta_2} \cdot \mathrm{adj}(M_{\beta_1})$ that some
implementations (including ours) use as an alternative to
computing $\theta = \beta_2 \cdot \overline{\beta_1} /
\mathrm{nrd}(I)$ as a quaternion does \emph{not} work for
cross-order pairs: it requires both $M_{\beta_i}$ to be
well-defined relative to a common torsion basis, which the
$\mathcal{O}_s, \mathcal{O}_t$ decompositions are not.  An
implementer who chooses this factoring will pass every
single-order KAT and hit the cross-order path as a hard
\textsc{Decompose}-returns-\texttt{None} error rather than a
quiet wrong-result; this is a comparatively kind failure mode,
but only because both bugs cancel.  An implementer who instead
follows \algo{3.13}'s natural reading and decomposes on
$\mathcal{O}_s, \mathcal{O}_t$ directly hits the same error.
Either way, the spec sends the implementer to the C reference
to discover the always-decompose-on-$\mathcal{O}_0$ convention.

% =====================================================================
\section{Sign error in \algo{4.9} line~11 aux-basis scaling}
\label{sec:alg49-line11-sign}

\finding{Affects \algo{4.9} (\textsc{Verify}), line~11.}{\high}

\algo{4.9} line~11 reads
\[
  P_{\mathrm{aux}}, Q_{\mathrm{aux}}
    \leftarrow [2^{f - e_{\mathrm{rsp}}' + 2}]\,P_{\mathrm{aux}},
               [2^{f - e_{\mathrm{rsp}}' + 2}]\,Q_{\mathrm{aux}}
\]
with a \emph{positive} sign on the $+2$.  After multiplying by
$2^{k}$ a point of order $2^{f}$ has order $2^{f - k}$, so
$k = f - e_{\mathrm{rsp}}' + 2$ leaves the aux basis at order
$2^{e_{\mathrm{rsp}}' - 2}$.  That is two bits too low for the
$(2,2)$-chain at line~26, which consumes
$e_{\mathrm{rsp}}' + 2$ bits of torsion (including the two
\texttt{HD\_extra\_torsion} bits the chain absorbs in its final
two steps).

The analogous line~13 reads
$[2^{f - e_{\mathrm{rsp}}' - r_{\mathrm{rsp}} - 2}]\,P_{\mathrm{chl}}$
with a \emph{negative} sign — consistent with the same chain
consuming $e_{\mathrm{rsp}}' + 2$ bits plus the additional
$r_{\mathrm{rsp}}$ bits the even-response isogeny eats.  Line~11
should match in sign and structure: $k = f - e_{\mathrm{rsp}}' - 2$,
landing the aux basis at order $2^{e_{\mathrm{rsp}}' + 2}$.

The C reference uses the negative-sign exponent (the
\texttt{f - pow\_dim2\_deg\_resp - 2} doubling count in
\texttt{challenge\_and\_aux\_basis\_verify}, equivalent to
$f - e_{\mathrm{rsp}}' - 2$).  An implementer who follows the
spec literally writes the positive-sign exponent, produces an
aux basis two bits short of the chain's torsion budget, and
line~26 fails with the all-zero-theta-null symptom shared by
several other bugs in this document.

\paragraph{Recommendation.}  Correct the sign in \algo{4.9}
line~11 to $f - e_{\mathrm{rsp}}' - 2$, matching line~13's
structure and the C reference.

% =====================================================================
\section{No parse-time canonicality on $\mathrm{chl}$'s top bits}
\label{sec:chl-canonicality}

\finding{Affects \algo{4.9} (\textsc{Verify}), line~2 (signature
parse).}{\low}

The challenge value $\mathrm{chl}$ in the signature is encoded
in $\lceil \texttt{SECURITY\_BITS} / 8 \rceil = 16$ bytes for
NIST-I, but the spec bounds the value by
$2^{e_{\mathrm{chl}}} = 2^{122}$.  The encoding allocates six
high bits (byte~15 bits 122--127) that the spec does not
constrain at parse time.  An honestly-produced signature has
those bits zero because the hash output is masked by
\texttt{CHALLENGE\_TOP\_MASK} before being written, but a
hostile signer can set them arbitrarily without violating any
spec-stated condition.

The C reference's \texttt{signature\_from\_bytes}
(\texttt{encode\_verification.c:191}) decodes the 16 chl bytes
verbatim without a top-bit check.  A signature whose chl bytes
have non-zero top bits parses successfully and proceeds into
the ladder
$[\mathrm{chl}]\,(P_{\mathrm{pk}} + [\mathrm{chl}]\,Q_{\mathrm{pk}})$
at line~9, where the extra bits steer the kernel to a different
$E_{\mathrm{chl}}$ than the honest signer intended.  The final
equality at line~30 compares the malformed stored chl against
the freshly-hashed $\mathrm{chl}'$ (whose top bits are masked);
the mismatch rejects the signature.

So the spec is functionally correct: malformed top bits surface
as $\mathrm{chl} \ne \mathrm{chl}'$ and verification fails.  But
the bad input runs the entire challenge ladder and the response
$(2,2)$-chain before being rejected.  A parse-time top-bit
check would reject in microseconds where the chain-then-compare
path costs milliseconds, which matters for DoS resistance under
spam.

\paragraph{Recommendation.}  Note in \algo{4.9} line~2 that
implementations may reject $\mathrm{chl} \ge 2^{e_{\mathrm{chl}}}$
at parse time without changing the semantics; doing so is a
defense-in-depth choice that does not alter the accept/reject
outcome on any signature.

% =====================================================================
\section{\algo{4.9} line 24: lower bound on $e'_{\mathrm{rsp}}$ for
  the $(2,2)$-chain entry}
\label{sec:spec-e-rsp-prime-lower-bound}

\finding{Affects \algo{4.9} (\textsc{Verify}, line 24) and the
$(2,2)$-chain machine \algo{8.47}
(\textsc{Isogeny22ChainWithTorsion}).}{\medium}

\algo{4.9} reads off $e'_{\mathrm{rsp}} = e_{\mathrm{rsp}} -
n_{\mathrm{bt}} - r_{\mathrm{rsp}}$ from the signature
(line 7) and uses it as the length of the verification
$(2,2)$-chain at line 24.  Line 7's check requires only
$e'_{\mathrm{rsp}} \geq 0$.

\paragraph{What goes wrong.}
\algo{8.47}'s extra-torsion mode runs a balanced strategy whose
indexing assumes \emph{at least one main step after the gluing
prep}.  Concretely, the chain needs the kernel to support
gluing $+\geq 1$ main step, i.e.\ $e'_{\mathrm{rsp}} \geq 2$ at
verification.  The spec text does not state this, and the only
guard at parse time
($e'_{\mathrm{rsp}} \geq 0$) admits $e'_{\mathrm{rsp}}
\in \{0, 1\}$.

For $e'_{\mathrm{rsp}} = 0$ the algorithm has an explicit
short-circuit branch (\algo{4.9} lines 21--23, ``skip the
$(2,2)$-isogeny''), so that case is handled.  The
$e'_{\mathrm{rsp}} = 1$ case has no such carve-out:
verification proceeds into the chain machine with chain length
$1$, and a straightforward implementation of \algo{8.47} either
indexes off the end of the strategy array, asserts that the
kernel ``has order $2^{e+2}$'' against an actual order of $2^3$,
or otherwise crashes.  An adversary need only flip
$n_{\mathrm{bt}}$ or $r_{\mathrm{rsp}}$ so that
$n_{\mathrm{bt}} + r_{\mathrm{rsp}} = e_{\mathrm{rsp}} - 1$ and
pick a $M_{\mathrm{chl}}$ that survives the $2^3$-isotropy guard;
verification crashes on every release build with a
non-conditional \texttt{assert!}.

\paragraph{Recommendation.}
\algo{4.9} line 7 (or an adjacent check) should require
$e'_{\mathrm{rsp}} \geq 2$ \emph{or} $e'_{\mathrm{rsp}} = 0$,
matching the two branches \algo{4.9} actually supports.
Equivalently: define $e'_{\mathrm{rsp}} = 1$ as an explicit
rejection.  The Wycheproof SQIsign suite includes a regression
vector for this case (\texttt{ShortChain} flag).

% =====================================================================
\section{\algo{8.47} extra-torsion=false mode: lower bound on
  $\mathrm{sui}.e$}
\label{sec:spec-isogeny-no-extra-torsion-bound}

\finding{Affects \algo{8.47}
(\textsc{Isogeny22ChainWithTorsion}, the
\texttt{extra\_torsion} = \texttt{false} variant), invoked from
the signing-side \textsc{IdealToIsogeny} outer chain
(\algo{3.13}).}{\medium}

\algo{8.47}'s \texttt{extra\_torsion} = \texttt{false} branch has
a different chain layout: gluing prep, then $e - 2$ main
$(2,2)$-steps, then a 4-isogeny tail, then a 2-isogeny tail
(producing the dual form for the splitter without padding the
kernel up to $2^{e+2}$).  This decomposition requires
$e \geq 4$~--- gluing plus one main step plus the two tails
exhausts the kernel only when there are four steps' worth of
torsion to consume.

The spec presents \algo{8.47} as a single algorithm parametrized
by \texttt{extra\_torsion}; the lower-bound asymmetry between
the two modes is not stated.

\paragraph{What goes wrong.}
\textsc{SuitableIdeals} (\algo{3.16}) selects an exponent
$\mathrm{sui}.e$ via the \textsc{FindUV} inner loop.  Honest
signing on NIST-I picks $\mathrm{sui}.e \approx f - 2 = 246$,
nowhere near the boundary.  A malformed signing key
(corrupted-on-disk, parser-accepted) can drive \textsc{FindUV}
into a tiny $\mathrm{sui}.e$~--- our fuzz harness for
parse-arbitrary-SK-then-sign surfaced such inputs in a
$15$-minute run.  The chain machine's
$\mathrm{sui}.e < 4$ then trips an
\texttt{assert!} or out-of-range index that the
release-mode build does not catch.  This is a DoS against a
defender's signer process whenever the secret key file is
tampered with.

\paragraph{Recommendation.}
State the lower bound explicitly in \algo{8.47}:
\texttt{extra\_torsion} = \texttt{true} requires
$e \geq 2$; \texttt{extra\_torsion} = \texttt{false} requires
$e \geq 4$.  Implementers should treat the bound as a guard
that returns a structured ``cannot run this chain'' error;
sign-side callers should consume that error as a chain failure
in their retry loop, not as a panic.

The same advice applies more generally to chain-machine
algorithms in the spec: every chain length~$e$ has a minimum
that depends on the machine's tail layout, and those minima
should be stated alongside the algorithm rather than implied by
the index ranges in the pseudocode.

\end{document}
